% Môn: Vật lý
% Lớp: 11
% Chương: Chương II. Sóng
% Bài: L11C2 Bài 12. Giao thoa sóng
% Số câu: 169
% App đọc file này trực tiếp — sửa rồi Commit trên GitHub, bấm Đồng bộ GitHub .tex

% L11C2 Bài 12. Giao thoa sóng

% ===== Câu 1 =====
% ID: L11C2B12-01-DS
% Mức: TH
\dangbt{Bài toán tổng hợp về giao thoa hai nguồn}
\begin{ex}
% ID: L11C2B12-01-DS
% Mức: TH
Xét các phát biểu sau trong dạng “Bài toán tổng hợp về giao thoa hai nguồn”.
\choiceTF
{\True Hai nguồn kết hợp là hai nguồn dao động — phát biểu đúng là: cùng phương, cùng tần số và có độ lệch pha không đổi.}
{Hai nguồn kết hợp là hai nguồn dao động — phát biểu cùng biên độ nhưng khác tần số là đúng.}
{\True Điều kiện để tại $M$ có cực đại giao thoa của hai nguồn cùng pha là — phát biểu đúng là: $d_2-d_1=k\lambda$.}
{Điều kiện để tại $M$ có cực đại giao thoa của hai nguồn cùng pha là — phát biểu $d_2+d_1=k\lambda$ là đúng.}
\loigiai{
\begin{itemchoice}
\item (Đúng) Hai nguồn kết hợp là hai nguồn dao động — phát biểu đúng là: cùng phương, cùng tần số và có độ lệch pha không đổi. (Vì): Hai nguồn kết hợp phải cùng phương, cùng tần số và có độ lệch pha không đổi. b) Sai: phương án nêu không phù hợp. c) Đúng: Với hai nguồn cùng pha, cực đại khi hiệu đường đi bằng số nguyên lần bước sóng. d) Sai: phương án nêu không phù hợp
\item (Sai) Hai nguồn kết hợp là hai nguồn dao động — phát biểu cùng biên độ nhưng khác tần số là đúng.
\item (Đúng) Điều kiện để tại $M$ có cực đại giao thoa của hai nguồn cùng pha là — phát biểu đúng là: $d_2-d_1=k\lambda$.
\item (Sai) Điều kiện để tại $M$ có cực đại giao thoa của hai nguồn cùng pha là — phát biểu $d_2+d_1=k\lambda$ là đúng.
\end{itemchoice}
}
\end{ex}

% ===== Câu 2 =====
% ID: L11C2B12-01-TL
% Mức: TH
\begin{ex}
% ID: L11C2B12-01-TL
% Mức: TH
Trong dạng “Bài toán tổng hợp về giao thoa hai nguồn”, Nêu điều kiện để hai nguồn sóng là hai nguồn kết hợp.
\choiceTF

\loigiai{
Hai nguồn dao động cùng phương, cùng tần số và có độ lệch pha không đổi theo thời gian.
}
\end{ex}

% ===== Câu 3 =====
% ID: L11C2B12-01-TLN
% Mức: TH
\begin{ex}
% ID: L11C2B12-01-TLN
% Mức: TH
Trong dạng “Bài toán tổng hợp về giao thoa hai nguồn”, Hai nguồn cùng pha, $\lambda=3$ cm. Tại $M$, hiệu đường đi là $12\,\text{cm}$. Xác định bậc cực đại.
\shortans{4}
\loigiai{
$k=(d_2-d_1)/\lambda=12/3=4$.
}
\end{ex}

% ===== Câu 4 =====
% ID: L11C2B12-01-TN
% Mức: NB
\begin{ex}
% ID: L11C2B12-01-TN
% Mức: NB
Trong dạng “Bài toán tổng hợp về giao thoa hai nguồn”, Hai nguồn kết hợp là hai nguồn dao động
\choice
{\True cùng phương, cùng tần số và có độ lệch pha không đổi}
{cùng biên độ nhưng khác tần số}
{cùng pha ở mọi thời điểm dù khác tần số}
{vuông pha và khác phương}
\loigiai{
Đáp án A. Hai nguồn kết hợp phải cùng phương, cùng tần số và có độ lệch pha không đổi.
}
\end{ex}

% ===== Câu 5 =====
% ID: L11C2B12-02-DS
% Mức: TH
\begin{ex}
% ID: L11C2B12-02-DS
% Mức: TH
Xét các phát biểu sau trong dạng “Bài toán tổng hợp về giao thoa hai nguồn”.
\choiceTF
{\True Điều kiện để tại $M$ có cực tiểu giao thoa của hai nguồn cùng pha là — phát biểu đúng là: $d_2-d_1=k\lambda$.}
{Điều kiện để tại $M$ có cực tiểu giao thoa của hai nguồn cùng pha là — phát biểu $d_2-d_1=(k+0,5)\lambda$ là đúng.}
{\True Hai nguồn cùng pha, $\lambda=2$ cm. Điểm $M$ có $d_2-d_1=6$ cm là — phát biểu đúng là: cực đại bậc 3.}
{Hai nguồn cùng pha, $\lambda=2$ cm. Điểm $M$ có $d_2-d_1=6$ cm là — phát biểu cực đại bậc 2 là đúng.}
\loigiai{
\begin{itemchoice}
\item (Đúng) Điều kiện để tại $M$ có cực tiểu giao thoa của hai nguồn cùng pha là — phát biểu đúng là: $d_2-d_1=k\lambda$. (Vì): Với hai nguồn cùng pha, cực tiểu khi hiệu đường đi bằng nửa nguyên lần bước sóng. b) Sai: phương án nêu không phù hợp. c) Đúng: Ta có $
\item (Sai) Điều kiện để tại $M$ có cực tiểu giao thoa của hai nguồn cùng pha là — phát biểu $d_2-d_1=(k+0,5)\lambda$ là đúng.
\item (Đúng) Hai nguồn cùng pha, $\lambda=2$ cm. Điểm $M$ có $d_2-d_1=6$ cm là — phát biểu đúng là: cực đại bậc 3.
\item (Sai) Hai nguồn cùng pha, $\lambda=2$ cm. Điểm $M$ có $d_2-d_1=6$ cm là — phát biểu cực đại bậc 2 là đúng. (Vì): _1)/\lambda=6/2=3$, nên $M$ thuộc cực đại bậc 3. d) Sai: phương án nêu không phù hợp
\end{itemchoice}
}
\end{ex}

% ===== Câu 6 =====
% ID: L11C2B12-02-TL
% Mức: TH
\begin{ex}
% ID: L11C2B12-02-TL
% Mức: TH
Trong dạng “Bài toán tổng hợp về giao thoa hai nguồn”, Giải thích vì sao ở cực đại giao thoa biên độ dao động lớn.
\choiceTF

\loigiai{
Tại cực đại, hai sóng đến cùng pha nên li độ tổng hợp được cộng đại số và biên độ đạt giá trị lớn.
}
\end{ex}

% ===== Câu 7 =====
% ID: L11C2B12-02-TLN
% Mức: TH
\begin{ex}
% ID: L11C2B12-02-TLN
% Mức: TH
Trong dạng “Bài toán tổng hợp về giao thoa hai nguồn”, Hai nguồn cùng pha, $\lambda=4$ cm. Hiệu đường đi tại cực tiểu thứ hai là bao nhiêu theo cm?
\shortans{6}
\loigiai{
Cực tiểu thứ hai ứng với $|d_2-d_1|=(1+0,5)\lambda=1,5\cdot4=6$ cm.
}
\end{ex}

% ===== Câu 8 =====
% ID: L11C2B12-02-TN
% Mức: NB
\begin{ex}
% ID: L11C2B12-02-TN
% Mức: NB
Trong dạng “Bài toán tổng hợp về giao thoa hai nguồn”, Điều kiện để tại $M$ có cực đại giao thoa của hai nguồn cùng pha là
\choice
{\True $d_2-d_1=k\lambda$}
{$d_2-d_1=(k+0,5)\lambda$}
{$d_2+d_1=k\lambda$}
{$d_2-d_1=\lambda/4$}
\loigiai{
Đáp án A. Với hai nguồn cùng pha, cực đại khi hiệu đường đi bằng số nguyên lần bước sóng.
}
\end{ex}

% ===== Câu 9 =====
% ID: L11C2B12-03-DS
% Mức: VD
\begin{ex}
% ID: L11C2B12-03-DS
% Mức: VD
Xét các phát biểu sau trong dạng “Bài toán tổng hợp về giao thoa hai nguồn”.
\choiceTF
{\True Hai nguồn cùng pha, $\lambda=4$ cm. Điểm $N$ có $d_2-d_1=6$ cm là — phát biểu đúng là: cực đại.}
{Hai nguồn cùng pha, $\lambda=4$ cm. Điểm $N$ có $d_2-d_1=6$ cm là — phát biểu cực tiểu là đúng.}
{\True Trên đoạn nối hai nguồn cùng pha, khoảng cách giữa hai cực đại liên tiếp bằng — phát biểu đúng là: $\lambda/2$.}
{Trên đoạn nối hai nguồn cùng pha, khoảng cách giữa hai cực đại liên tiếp bằng — phát biểu $2\lambda$ là đúng.}
\loigiai{
\begin{itemchoice}
\item (Đúng) Hai nguồn cùng pha, $\lambda=4$ cm. Điểm $N$ có $d_2-d_1=6$ cm là — phát biểu đúng là: cực đại. (Vì): $6/4=1,5$, là nửa nguyên nên $N$ thuộc cực tiểu. b) Sai: phương án nêu không phù hợp. c) Đúng: Trên đoạn nối hai nguồn, hiệu đường đi đổi hai lần theo li độ nên hai cực đại liên tiếp cách nhau $\lambda/2$. d) Sai: phương án nêu không phù hợp
\item (Sai) Hai nguồn cùng pha, $\lambda=4$ cm. Điểm $N$ có $d_2-d_1=6$ cm là — phát biểu cực tiểu là đúng.
\item (Đúng) Trên đoạn nối hai nguồn cùng pha, khoảng cách giữa hai cực đại liên tiếp bằng — phát biểu đúng là: $\lambda/2$.
\item (Sai) Trên đoạn nối hai nguồn cùng pha, khoảng cách giữa hai cực đại liên tiếp bằng — phát biểu $2\lambda$ là đúng.
\end{itemchoice}
}
\end{ex}

% ===== Câu 10 =====
% ID: L11C2B12-03-TL
% Mức: VD
\begin{ex}
% ID: L11C2B12-03-TL
% Mức: VD
Trong dạng “Bài toán tổng hợp về giao thoa hai nguồn”, Hai nguồn cùng pha có $\lambda=2$ cm. Xét điểm $M$ với $d_1=10$ cm, $d_2=15$ cm. Xác định trạng thái giao thoa tại M.
\choiceTF

\loigiai{
Hiệu đường đi $\Delta d=5\,\mathrm{cm}=2,5\lambda$, nên hai sóng ngược pha và $M$ thuộc cực tiểu.
}
\end{ex}

% ===== Câu 11 =====
% ID: L11C2B12-03-TLN
% Mức: VD
\begin{ex}
% ID: L11C2B12-03-TLN
% Mức: VD
Trong dạng “Bài toán tổng hợp về giao thoa hai nguồn”, Trên đoạn nối hai nguồn, hai cực đại liên tiếp cách nhau $2,5\,\text{cm}$. Tính bước sóng theo cm.
\shortans{5}
\loigiai{
Khoảng cách hai cực đại liên tiếp trên đoạn nối nguồn bằng $\lambda/2$, nên $\lambda=5$ cm.
}
\end{ex}

% ===== Câu 12 =====
% ID: L11C2B12-03-TN
% Mức: TH
\begin{ex}
% ID: L11C2B12-03-TN
% Mức: TH
Trong dạng “Bài toán tổng hợp về giao thoa hai nguồn”, Điều kiện để tại $M$ có cực tiểu giao thoa của hai nguồn cùng pha là
\choice
{$d_2-d_1=k\lambda$}
{\True $d_2-d_1=(k+0,5)\lambda$}
{$d_2+d_1=(k+0,5)\lambda$}
{$d_2-d_1=2k\lambda$}
\loigiai{
Đáp án B. Với hai nguồn cùng pha, cực tiểu khi hiệu đường đi bằng nửa nguyên lần bước sóng.
}
\end{ex}

% ===== Câu 13 =====
% ID: L11C2B12-04-DS
% Mức: VD
\begin{ex}
% ID: L11C2B12-04-DS
% Mức: VD
Xét các phát biểu sau trong dạng “Bài toán tổng hợp về giao thoa hai nguồn”.
\choiceTF
{\True Hiện tượng giao thoa là bằng chứng rõ ràng của — phát biểu đúng là: tính chất sóng.}
{Hiện tượng giao thoa là bằng chứng rõ ràng của — phát biểu tính chất hạt là đúng.}
{\True Hai nguồn cùng pha cách nhau $10\,\text{cm}$, $\lambda=2$ cm. Số cực đại trên đoạn nối hai nguồn, kể cả hai nguồn nếu thỏa điều kiện, là — phát biểu đúng là: 9.}
{Hai nguồn cùng pha cách nhau $10\,\text{cm}$, $\lambda=2$ cm. Số cực đại trên đoạn nối hai nguồn, kể cả hai nguồn nếu thỏa điều kiện, là — phát biểu 11 là đúng.}
\loigiai{
\begin{itemchoice}
\item (Đúng) Hiện tượng giao thoa là bằng chứng rõ ràng của — phát biểu đúng là: tính chất sóng. (Vì): Giao thoa là hiện tượng đặc trưng của sóng. b) Sai: phương án nêu không phù hợp. c) Đúng: Các bậc nguyên thỏa $|k|\le 5$, nên có 11 giá trị từ -5 đến 5. d) Sai: phương án nêu không phù hợp
\item (Sai) Hiện tượng giao thoa là bằng chứng rõ ràng của — phát biểu tính chất hạt là đúng.
\item (Đúng) Hai nguồn cùng pha cách nhau $10\,\text{cm}$, $\lambda=2$ cm. Số cực đại trên đoạn nối hai nguồn, kể cả hai nguồn nếu thỏa điều kiện, là — phát biểu đúng là: 9.
\item (Sai) Hai nguồn cùng pha cách nhau $10\,\text{cm}$, $\lambda=2$ cm. Số cực đại trên đoạn nối hai nguồn, kể cả hai nguồn nếu thỏa điều kiện, là — phát biểu 11 là đúng.
\end{itemchoice}
}
\end{ex}

% ===== Câu 14 =====
% ID: L11C2B12-04-TL
% Mức: VD
\begin{ex}
% ID: L11C2B12-04-TL
% Mức: VD
Trong dạng “Bài toán tổng hợp về giao thoa hai nguồn”, Trình bày cách xác định bước sóng từ ảnh các đường cực đại trên đoạn nối hai nguồn.
\choiceTF

\loigiai{
Đo khoảng cách $i$ giữa hai cực đại liên tiếp trên đoạn nối nguồn; do $i=\lambda/2$ nên tính $\lambda=2i$.
}
\end{ex}

% ===== Câu 15 =====
% ID: L11C2B12-04-TLN
% Mức: VD
\begin{ex}
% ID: L11C2B12-04-TLN
% Mức: VD
Trong dạng “Bài toán tổng hợp về giao thoa hai nguồn”, Sóng có tốc độ $40\,\text{cm}$ /s và tần số $10\,\text{Hz}$. Tính bước sóng theo cm.
\shortans{4}
\loigiai{
$\lambda=v/f=40/10=4$ cm.
}
\end{ex}

% ===== Câu 16 =====
% ID: L11C2B12-04-TN
% Mức: TH
\begin{ex}
% ID: L11C2B12-04-TN
% Mức: TH
Trong dạng “Bài toán tổng hợp về giao thoa hai nguồn”, Hai nguồn cùng pha, $\lambda=2$ cm. Điểm $M$ có $d_2-d_1=6$ cm là
\choice
{\True cực đại bậc 3}
{cực tiểu thứ 3}
{cực đại bậc 2}
{không giao thoa}
\loigiai{
Đáp án A. Ta có $(d_2-d_1)/\lambda=6/2=3$, nên $M$ thuộc cực đại bậc 3.
}
\end{ex}

% ===== Câu 17 =====
% ID: L11C2B12-05-DS
% Mức: VDC
\begin{ex}
% ID: L11C2B12-05-DS
% Mức: VDC
Xét các phát biểu sau trong dạng “Bài toán tổng hợp về giao thoa hai nguồn”.
\choiceTF
{\True Tại điểm có cực đại giao thoa, hai sóng thành phần đến đó — phát biểu đúng là: cùng pha.}
{Tại điểm có cực đại giao thoa, hai sóng thành phần đến đó — phát biểu ngược pha là đúng.}
{\True Nếu tần số nguồn tăng gấp đôi còn tốc độ truyền sóng không đổi thì bước sóng — phát biểu đúng là: giảm một nửa.}
{Nếu tần số nguồn tăng gấp đôi còn tốc độ truyền sóng không đổi thì bước sóng — phát biểu không đổi là đúng.}
\loigiai{
\begin{itemchoice}
\item (Đúng) Tại điểm có cực đại giao thoa, hai sóng thành phần đến đó — phát biểu đúng là: cùng pha. (Vì): Cực đại xuất hiện khi hai dao động thành phần tăng cường nhau, tức cùng pha. b) Sai: phương án nêu không phù hợp. c) Đúng: Từ $\lambda=v/f$, khi $f$ tăng gấp đôi thì $\lambda$ giảm một nửa. d) Sai: phương án nêu không phù hợp
\item (Sai) Tại điểm có cực đại giao thoa, hai sóng thành phần đến đó — phát biểu ngược pha là đúng.
\item (Đúng) Nếu tần số nguồn tăng gấp đôi còn tốc độ truyền sóng không đổi thì bước sóng — phát biểu đúng là: giảm một nửa.
\item (Sai) Nếu tần số nguồn tăng gấp đôi còn tốc độ truyền sóng không đổi thì bước sóng — phát biểu không đổi là đúng.
\end{itemchoice}
}
\end{ex}

% ===== Câu 18 =====
% ID: L11C2B12-05-TL
% Mức: VD
\begin{ex}
% ID: L11C2B12-05-TL
% Mức: VD
Trong dạng “Bài toán tổng hợp về giao thoa hai nguồn”, Hai nguồn cùng pha cách nhau $16\,\text{cm}$, bước sóng $4\,\text{cm}$. Xác định số cực đại trên đoạn nối hai nguồn.
\choiceTF

\loigiai{
Điều kiện $|k|\le16/4=4$. Có các giá trị $k=-4,-3,\ldots,4$, tổng cộng 9 cực đại.
}
\end{ex}

% ===== Câu 19 =====
% ID: L11C2B12-05-TLN
% Mức: VD
\begin{ex}
% ID: L11C2B12-05-TLN
% Mức: VD
Trong dạng “Bài toán tổng hợp về giao thoa hai nguồn”, Hai nguồn cùng pha cách nhau $12\,\text{cm}$, $\lambda=3$ cm. Giá trị bậc cực đại lớn nhất trên đoạn nối nguồn là bao nhiêu?
\shortans{4}
\loigiai{
$|k|\le a/\lambda=12/3=4$, nên bậc lớn nhất là 4.
}
\end{ex}

% ===== Câu 20 =====
% ID: L11C2B12-05-TN
% Mức: TH
\begin{ex}
% ID: L11C2B12-05-TN
% Mức: TH
Trong dạng “Bài toán tổng hợp về giao thoa hai nguồn”, Hai nguồn cùng pha, $\lambda=4$ cm. Điểm $N$ có $d_2-d_1=6$ cm là
\choice
{cực đại}
{\True cực tiểu}
{điểm đứng yên do không có sóng tới}
{không xác định}
\loigiai{
Đáp án B. $6/4=1,5$, là nửa nguyên nên $N$ thuộc cực tiểu.
}
\end{ex}

% ===== Câu 21 =====
% ID: L11C2B12-06-DS
% Mức: TH
\dangbt{Nhận biết hiện tượng và điều kiện giao thoa sóng}
\begin{ex}
% ID: L11C2B12-06-DS
% Mức: TH
Xét các phát biểu sau trong dạng “Nhận biết hiện tượng và điều kiện giao thoa sóng”.
\choiceTF
{\True Hai nguồn kết hợp là hai nguồn dao động — phát biểu đúng là: cùng phương, cùng tần số và có độ lệch pha không đổi.}
{Hai nguồn kết hợp là hai nguồn dao động — phát biểu cùng biên độ nhưng khác tần số là đúng.}
{\True Điều kiện để tại $M$ có cực đại giao thoa của hai nguồn cùng pha là — phát biểu đúng là: $d_2-d_1=k\lambda$.}
{Điều kiện để tại $M$ có cực đại giao thoa của hai nguồn cùng pha là — phát biểu $d_2+d_1=k\lambda$ là đúng.}
\loigiai{
\begin{itemchoice}
\item (Đúng) Hai nguồn kết hợp là hai nguồn dao động — phát biểu đúng là: cùng phương, cùng tần số và có độ lệch pha không đổi. (Vì): Hai nguồn kết hợp phải cùng phương, cùng tần số và có độ lệch pha không đổi. b) Sai: phương án nêu không phù hợp. c) Đúng: Với hai nguồn cùng pha, cực đại khi hiệu đường đi bằng số nguyên lần bước sóng. d) Sai: phương án nêu không phù hợp
\item (Sai) Hai nguồn kết hợp là hai nguồn dao động — phát biểu cùng biên độ nhưng khác tần số là đúng.
\item (Đúng) Điều kiện để tại $M$ có cực đại giao thoa của hai nguồn cùng pha là — phát biểu đúng là: $d_2-d_1=k\lambda$.
\item (Sai) Điều kiện để tại $M$ có cực đại giao thoa của hai nguồn cùng pha là — phát biểu $d_2+d_1=k\lambda$ là đúng.
\end{itemchoice}
}
\end{ex}

% ===== Câu 22 =====
% ID: L11C2B12-06-TL
% Mức: VDC
\dangbt{Bài toán tổng hợp về giao thoa hai nguồn}
\begin{ex}
% ID: L11C2B12-06-TL
% Mức: VDC
Trong dạng “Bài toán tổng hợp về giao thoa hai nguồn”, Phân biệt cực đại và cực tiểu giao thoa của hai nguồn cùng pha bằng hiệu đường đi.
\choiceTF

\loigiai{
Cực đại: $\Delta d=k\lambda$. Cực tiểu: $\Delta d=(k+1/2)\lambda$, với $k$ nguyên.
}
\end{ex}

% ===== Câu 23 =====
% ID: L11C2B12-06-TLN
% Mức: VDC
\begin{ex}
% ID: L11C2B12-06-TLN
% Mức: VDC
Trong dạng “Bài toán tổng hợp về giao thoa hai nguồn”, Tại $M$, $d_1=18$ cm, $d_2=25$ cm. Tính độ lớn hiệu đường đi theo cm.
\shortans{7}
\loigiai{
$\Delta d=|d_2-d_1|=|25-18|=7$ cm.
}
\end{ex}

% ===== Câu 24 =====
% ID: L11C2B12-06-TN
% Mức: VD
\begin{ex}
% ID: L11C2B12-06-TN
% Mức: VD
Trong dạng “Bài toán tổng hợp về giao thoa hai nguồn”, Trên đoạn nối hai nguồn cùng pha, khoảng cách giữa hai cực đại liên tiếp bằng
\choice
{\True $\lambda/2$}
{$\lambda$}
{$2\lambda$}
{$\lambda/4$}
\loigiai{
Đáp án A. Trên đoạn nối hai nguồn, hiệu đường đi đổi hai lần theo li độ nên hai cực đại liên tiếp cách nhau $\lambda/2$.
}
\end{ex}

% ===== Câu 25 =====
% ID: L11C2B12-07-DS
% Mức: TH
\dangbt{Nhận biết hiện tượng và điều kiện giao thoa sóng}
\begin{ex}
% ID: L11C2B12-07-DS
% Mức: TH
Xét các phát biểu sau trong dạng “Nhận biết hiện tượng và điều kiện giao thoa sóng”.
\choiceTF
{\True Điều kiện để tại $M$ có cực tiểu giao thoa của hai nguồn cùng pha là — phát biểu đúng là: $d_2-d_1=k\lambda$.}
{Điều kiện để tại $M$ có cực tiểu giao thoa của hai nguồn cùng pha là — phát biểu $d_2-d_1=(k+0,5)\lambda$ là đúng.}
{\True Hai nguồn cùng pha, $\lambda=2$ cm. Điểm $M$ có $d_2-d_1=6$ cm là — phát biểu đúng là: cực đại bậc 3.}
{Hai nguồn cùng pha, $\lambda=2$ cm. Điểm $M$ có $d_2-d_1=6$ cm là — phát biểu cực đại bậc 2 là đúng.}
\loigiai{
\begin{itemchoice}
\item (Đúng) Điều kiện để tại $M$ có cực tiểu giao thoa của hai nguồn cùng pha là — phát biểu đúng là: $d_2-d_1=k\lambda$. (Vì): Với hai nguồn cùng pha, cực tiểu khi hiệu đường đi bằng nửa nguyên lần bước sóng. b) Sai: phương án nêu không phù hợp. c) Đúng: Ta có $
\item (Sai) Điều kiện để tại $M$ có cực tiểu giao thoa của hai nguồn cùng pha là — phát biểu $d_2-d_1=(k+0,5)\lambda$ là đúng.
\item (Đúng) Hai nguồn cùng pha, $\lambda=2$ cm. Điểm $M$ có $d_2-d_1=6$ cm là — phát biểu đúng là: cực đại bậc 3.
\item (Sai) Hai nguồn cùng pha, $\lambda=2$ cm. Điểm $M$ có $d_2-d_1=6$ cm là — phát biểu cực đại bậc 2 là đúng. (Vì): _1)/\lambda=6/2=3$, nên $M$ thuộc cực đại bậc 3. d) Sai: phương án nêu không phù hợp
\end{itemchoice}
}
\end{ex}

% ===== Câu 26 =====
% ID: L11C2B12-07-TL
% Mức: TH
\dangbt{Chưa có dạng}
\begin{ex}
% ID: L11C2B12-07-TL
% Mức: TH
Trong thí nghiệm Young về giao thoa với ánh sáng đơn sắc có bước sóng $\lambda,$ khoảng cách giữa hai khe là $0,15 \mathrm{mm}$, khoảng cách giữa mặt phẳng chứa hai khe và màn quan sát là $1\,\text{m}$. Hai điểm $M$ và $N$ trên màn quan sát đối xứng nhau qua vân sáng trung tâm. Trên đoạn MN có 11 vân sáng, tại $M$ và $N$ là hai vân sáng. Biết khoảng cách MN là $30 \mathrm{mm}$. Tính bước sóng của ánh sáng dùng trong thí nghiệm này. Câu 11*. Trong thí nghiệm Young về giao thoa ánh sáng, nguồn $S$ phát ánh sáng đơn sắc có bước sóng $\lambda$. Màn quan sát cách hai khe một khoảng không đổi $D$, khoảng cách giữa hai khe S_1 $S_2 = a có thể thay đổi (nhưng S_1$, S_2 luôn cách đều S). Xét điểm $P$ trên màn quan sát, lúc đầu là vân sáng bậc 4, nếu lần lượt giảm hoặc tăng khoảng cách S_1 S_2 một lượng $\Deltaa$ thì tại đó là vân sáng bậc k và 3k. Nếu tăng khoảng cách S_1 S_2 một lượng 2 $\Delta$ a thì tại đó là vân sáng hay vân tối, bậc hoặc thứ bao nhiêu? Câu 12*. Thực hiện thí nghiệm Young về giao thoa với ánh sáng, khoảng cách giữa hai khe là $2 \mathrm{mm}$, khoảng cách giữa mặt phẳng chứa hai khe và màn quan sát là $2\,\text{m}$. Người ta chiếu sáng hai khe bằng ánh sáng trắng có bước sóng nằm trong khoảng từ 0,40 $\mum$ đến 0,76 $\mum$. Hỏi tại điểm $M$ trên màn ảnh cách vân sáng trung tâm $3,3 \mathrm{mm}$ sẽ cho vân tối có bước sóng ngắn nhất bằng bao nhiêu? BÀl 13. SÓNG DỪNG $\nguon{ly11 kntt.pdf}$
\choiceTF

\loigiai{
$\text{i} = \dfrac{\text{MN}}{11 - 1} = \dfrac{30 \cdot 10^{-3}}{11 - 1} = 3\cdot 10^{-3}\text{m}\left.\Rightarrow\lambda=\dfrac{\text{ai}}{\text{D}}=\dfrac{0,15 \cdot 10^{-3} \cdot 3 \cdot 10^{-3}}{1}= 4,5\cdot10^{-7}\text{m}.\right.$
}
\end{ex}

% ===== Câu 27 =====
% ID: L11C2B12-07-TLN
% Mức: TH
\dangbt{Nhận biết hiện tượng và điều kiện giao thoa sóng}
\begin{ex}
% ID: L11C2B12-07-TLN
% Mức: TH
Trong dạng “Nhận biết hiện tượng và điều kiện giao thoa sóng”, Hai nguồn cùng pha, $\lambda=3$ cm. Tại $M$, hiệu đường đi là $12\,\text{cm}$. Xác định bậc cực đại.
\shortans{4}
\loigiai{
$k=(d_2-d_1)/\lambda=12/3=4$.
}
\end{ex}

% ===== Câu 28 =====
% ID: L11C2B12-07-TN
% Mức: VD
\dangbt{Bài toán tổng hợp về giao thoa hai nguồn}
\begin{ex}
% ID: L11C2B12-07-TN
% Mức: VD
Trong dạng “Bài toán tổng hợp về giao thoa hai nguồn”, Hiện tượng giao thoa là bằng chứng rõ ràng của
\choice
{\True tính chất sóng}
{tính chất hạt}
{sự bảo toàn khối lượng}
{chuyển động thẳng đều}
\loigiai{
Đáp án A. Giao thoa là hiện tượng đặc trưng của sóng.
}
\end{ex}

% ===== Câu 29 =====
% ID: L11C2B12-08-DS
% Mức: VD
\dangbt{Nhận biết hiện tượng và điều kiện giao thoa sóng}
\begin{ex}
% ID: L11C2B12-08-DS
% Mức: VD
Xét các phát biểu sau trong dạng “Nhận biết hiện tượng và điều kiện giao thoa sóng”.
\choiceTF
{\True Hai nguồn cùng pha, $\lambda=4$ cm. Điểm $N$ có $d_2-d_1=6$ cm là — phát biểu đúng là: cực đại.}
{Hai nguồn cùng pha, $\lambda=4$ cm. Điểm $N$ có $d_2-d_1=6$ cm là — phát biểu cực tiểu là đúng.}
{\True Trên đoạn nối hai nguồn cùng pha, khoảng cách giữa hai cực đại liên tiếp bằng — phát biểu đúng là: $\lambda/2$.}
{Trên đoạn nối hai nguồn cùng pha, khoảng cách giữa hai cực đại liên tiếp bằng — phát biểu $2\lambda$ là đúng.}
\loigiai{
\begin{itemchoice}
\item (Đúng) Hai nguồn cùng pha, $\lambda=4$ cm. Điểm $N$ có $d_2-d_1=6$ cm là — phát biểu đúng là: cực đại. (Vì): $6/4=1,5$, là nửa nguyên nên $N$ thuộc cực tiểu. b) Sai: phương án nêu không phù hợp. c) Đúng: Trên đoạn nối hai nguồn, hiệu đường đi đổi hai lần theo li độ nên hai cực đại liên tiếp cách nhau $\lambda/2$. d) Sai: phương án nêu không phù hợp
\item (Sai) Hai nguồn cùng pha, $\lambda=4$ cm. Điểm $N$ có $d_2-d_1=6$ cm là — phát biểu cực tiểu là đúng.
\item (Đúng) Trên đoạn nối hai nguồn cùng pha, khoảng cách giữa hai cực đại liên tiếp bằng — phát biểu đúng là: $\lambda/2$.
\item (Sai) Trên đoạn nối hai nguồn cùng pha, khoảng cách giữa hai cực đại liên tiếp bằng — phát biểu $2\lambda$ là đúng.
\end{itemchoice}
}
\end{ex}

% ===== Câu 30 =====
% ID: L11C2B12-08-TL
% Mức: TH
\begin{ex}
% ID: L11C2B12-08-TL
% Mức: TH
Trong dạng “Nhận biết hiện tượng và điều kiện giao thoa sóng”, Nêu điều kiện để hai nguồn sóng là hai nguồn kết hợp.
\choiceTF

\loigiai{
Hai nguồn dao động cùng phương, cùng tần số và có độ lệch pha không đổi theo thời gian.
}
\end{ex}

% ===== Câu 31 =====
% ID: L11C2B12-08-TLN
% Mức: TH
\begin{ex}
% ID: L11C2B12-08-TLN
% Mức: TH
Trong dạng “Nhận biết hiện tượng và điều kiện giao thoa sóng”, Hai nguồn cùng pha, $\lambda=4$ cm. Hiệu đường đi tại cực tiểu thứ hai là bao nhiêu theo cm?
\shortans{6}
\loigiai{
Cực tiểu thứ hai ứng với $|d_2-d_1|=(1+0,5)\lambda=1,5\cdot4=6$ cm.
}
\end{ex}

% ===== Câu 32 =====
% ID: L11C2B12-08-TN
% Mức: VD
\dangbt{Bài toán tổng hợp về giao thoa hai nguồn}
\begin{ex}
% ID: L11C2B12-08-TN
% Mức: VD
Trong dạng “Bài toán tổng hợp về giao thoa hai nguồn”, Hai nguồn cùng pha cách nhau $10\,\text{cm}$, $\lambda=2$ cm. Số cực đại trên đoạn nối hai nguồn, kể cả hai nguồn nếu thỏa điều kiện, là
\choice
{9}
{10}
{\True 11}
{12}
\loigiai{
Đáp án C. Các bậc nguyên thỏa $|k|\le 5$, nên có 11 giá trị từ -5 đến 5.
}
\end{ex}

% ===== Câu 33 =====
% ID: L11C2B12-09-DS
% Mức: VD
\dangbt{Nhận biết hiện tượng và điều kiện giao thoa sóng}
\begin{ex}
% ID: L11C2B12-09-DS
% Mức: VD
Xét các phát biểu sau trong dạng “Nhận biết hiện tượng và điều kiện giao thoa sóng”.
\choiceTF
{\True Hiện tượng giao thoa là bằng chứng rõ ràng của — phát biểu đúng là: tính chất sóng.}
{Hiện tượng giao thoa là bằng chứng rõ ràng của — phát biểu tính chất hạt là đúng.}
{\True Hai nguồn cùng pha cách nhau $10\,\text{cm}$, $\lambda=2$ cm. Số cực đại trên đoạn nối hai nguồn, kể cả hai nguồn nếu thỏa điều kiện, là — phát biểu đúng là: 9.}
{Hai nguồn cùng pha cách nhau $10\,\text{cm}$, $\lambda=2$ cm. Số cực đại trên đoạn nối hai nguồn, kể cả hai nguồn nếu thỏa điều kiện, là — phát biểu 11 là đúng.}
\loigiai{
\begin{itemchoice}
\item (Đúng) Hiện tượng giao thoa là bằng chứng rõ ràng của — phát biểu đúng là: tính chất sóng. (Vì): Giao thoa là hiện tượng đặc trưng của sóng. b) Sai: phương án nêu không phù hợp. c) Đúng: Các bậc nguyên thỏa $|k|\le 5$, nên có 11 giá trị từ -5 đến 5. d) Sai: phương án nêu không phù hợp
\item (Sai) Hiện tượng giao thoa là bằng chứng rõ ràng của — phát biểu tính chất hạt là đúng.
\item (Đúng) Hai nguồn cùng pha cách nhau $10\,\text{cm}$, $\lambda=2$ cm. Số cực đại trên đoạn nối hai nguồn, kể cả hai nguồn nếu thỏa điều kiện, là — phát biểu đúng là: 9.
\item (Sai) Hai nguồn cùng pha cách nhau $10\,\text{cm}$, $\lambda=2$ cm. Số cực đại trên đoạn nối hai nguồn, kể cả hai nguồn nếu thỏa điều kiện, là — phát biểu 11 là đúng.
\end{itemchoice}
}
\end{ex}

% ===== Câu 34 =====
% ID: L11C2B12-09-TL
% Mức: TH
\begin{ex}
% ID: L11C2B12-09-TL
% Mức: TH
Trong dạng “Nhận biết hiện tượng và điều kiện giao thoa sóng”, Giải thích vì sao ở cực đại giao thoa biên độ dao động lớn.
\choiceTF

\loigiai{
Tại cực đại, hai sóng đến cùng pha nên li độ tổng hợp được cộng đại số và biên độ đạt giá trị lớn.
}
\end{ex}

% ===== Câu 35 =====
% ID: L11C2B12-09-TLN
% Mức: VD
\begin{ex}
% ID: L11C2B12-09-TLN
% Mức: VD
Trong dạng “Nhận biết hiện tượng và điều kiện giao thoa sóng”, Trên đoạn nối hai nguồn, hai cực đại liên tiếp cách nhau $2,5\,\text{cm}$. Tính bước sóng theo cm.
\shortans{5}
\loigiai{
Khoảng cách hai cực đại liên tiếp trên đoạn nối nguồn bằng $\lambda/2$, nên $\lambda=5$ cm.
}
\end{ex}

% ===== Câu 36 =====
% ID: L11C2B12-09-TN
% Mức: VDC
\dangbt{Bài toán tổng hợp về giao thoa hai nguồn}
\begin{ex}
% ID: L11C2B12-09-TN
% Mức: VDC
Trong dạng “Bài toán tổng hợp về giao thoa hai nguồn”, Tại điểm có cực đại giao thoa, hai sóng thành phần đến đó
\choice
{\True cùng pha}
{ngược pha}
{vuông pha}
{khác tần số}
\loigiai{
Đáp án A. Cực đại xuất hiện khi hai dao động thành phần tăng cường nhau, tức cùng pha.
}
\end{ex}

% ===== Câu 37 =====
% ID: L11C2B12-10-DS
% Mức: VDC
\dangbt{Nhận biết hiện tượng và điều kiện giao thoa sóng}
\begin{ex}
% ID: L11C2B12-10-DS
% Mức: VDC
Xét các phát biểu sau trong dạng “Nhận biết hiện tượng và điều kiện giao thoa sóng”.
\choiceTF
{\True Tại điểm có cực đại giao thoa, hai sóng thành phần đến đó — phát biểu đúng là: cùng pha.}
{Tại điểm có cực đại giao thoa, hai sóng thành phần đến đó — phát biểu ngược pha là đúng.}
{\True Nếu tần số nguồn tăng gấp đôi còn tốc độ truyền sóng không đổi thì bước sóng — phát biểu đúng là: giảm một nửa.}
{Nếu tần số nguồn tăng gấp đôi còn tốc độ truyền sóng không đổi thì bước sóng — phát biểu không đổi là đúng.}
\loigiai{
\begin{itemchoice}
\item (Đúng) Tại điểm có cực đại giao thoa, hai sóng thành phần đến đó — phát biểu đúng là: cùng pha. (Vì): Cực đại xuất hiện khi hai dao động thành phần tăng cường nhau, tức cùng pha. b) Sai: phương án nêu không phù hợp. c) Đúng: Từ $\lambda=v/f$, khi $f$ tăng gấp đôi thì $\lambda$ giảm một nửa. d) Sai: phương án nêu không phù hợp
\item (Sai) Tại điểm có cực đại giao thoa, hai sóng thành phần đến đó — phát biểu ngược pha là đúng.
\item (Đúng) Nếu tần số nguồn tăng gấp đôi còn tốc độ truyền sóng không đổi thì bước sóng — phát biểu đúng là: giảm một nửa.
\item (Sai) Nếu tần số nguồn tăng gấp đôi còn tốc độ truyền sóng không đổi thì bước sóng — phát biểu không đổi là đúng.
\end{itemchoice}
}
\end{ex}

% ===== Câu 38 =====
% ID: L11C2B12-10-TL
% Mức: VD
\begin{ex}
% ID: L11C2B12-10-TL
% Mức: VD
Trong dạng “Nhận biết hiện tượng và điều kiện giao thoa sóng”, Hai nguồn cùng pha có $\lambda=2$ cm. Xét điểm $M$ với $d_1=10$ cm, $d_2=15$ cm. Xác định trạng thái giao thoa tại M.
\choiceTF

\loigiai{
Hiệu đường đi $\Delta d=5\,\mathrm{cm}=2,5\lambda$, nên hai sóng ngược pha và $M$ thuộc cực tiểu.
}
\end{ex}

% ===== Câu 39 =====
% ID: L11C2B12-10-TLN
% Mức: VD
\begin{ex}
% ID: L11C2B12-10-TLN
% Mức: VD
Trong dạng “Nhận biết hiện tượng và điều kiện giao thoa sóng”, Sóng có tốc độ $40\,\text{cm}$ /s và tần số $10\,\text{Hz}$. Tính bước sóng theo cm.
\shortans{4}
\loigiai{
$\lambda=v/f=40/10=4$ cm.
}
\end{ex}

% ===== Câu 40 =====
% ID: L11C2B12-10-TN
% Mức: VDC
\dangbt{Bài toán tổng hợp về giao thoa hai nguồn}
\begin{ex}
% ID: L11C2B12-10-TN
% Mức: VDC
Trong dạng “Bài toán tổng hợp về giao thoa hai nguồn”, Nếu tần số nguồn tăng gấp đôi còn tốc độ truyền sóng không đổi thì bước sóng
\choice
{\True giảm một nửa}
{tăng gấp đôi}
{không đổi}
{tăng bốn lần}
\loigiai{
Đáp án A. Từ $\lambda=v/f$, khi $f$ tăng gấp đôi thì $\lambda$ giảm một nửa.
}
\end{ex}

% ===== Câu 41 =====
% ID: L11C2B12-11-DS
% Mức: TH
\dangbt{Tính hiệu đường đi và bậc giao thoa}
\begin{ex}
% ID: L11C2B12-11-DS
% Mức: TH
Xét các phát biểu sau trong dạng “Tính hiệu đường đi và bậc giao thoa”.
\choiceTF
{\True Hai nguồn kết hợp là hai nguồn dao động — phát biểu đúng là: cùng phương, cùng tần số và có độ lệch pha không đổi.}
{Hai nguồn kết hợp là hai nguồn dao động — phát biểu cùng biên độ nhưng khác tần số là đúng.}
{\True Điều kiện để tại $M$ có cực đại giao thoa của hai nguồn cùng pha là — phát biểu đúng là: $d_2-d_1=k\lambda$.}
{Điều kiện để tại $M$ có cực đại giao thoa của hai nguồn cùng pha là — phát biểu $d_2+d_1=k\lambda$ là đúng.}
\loigiai{
\begin{itemchoice}
\item (Đúng) Hai nguồn kết hợp là hai nguồn dao động — phát biểu đúng là: cùng phương, cùng tần số và có độ lệch pha không đổi. (Vì): Hai nguồn kết hợp phải cùng phương, cùng tần số và có độ lệch pha không đổi. b) Sai: phương án nêu không phù hợp. c) Đúng: Với hai nguồn cùng pha, cực đại khi hiệu đường đi bằng số nguyên lần bước sóng. d) Sai: phương án nêu không phù hợp
\item (Sai) Hai nguồn kết hợp là hai nguồn dao động — phát biểu cùng biên độ nhưng khác tần số là đúng.
\item (Đúng) Điều kiện để tại $M$ có cực đại giao thoa của hai nguồn cùng pha là — phát biểu đúng là: $d_2-d_1=k\lambda$.
\item (Sai) Điều kiện để tại $M$ có cực đại giao thoa của hai nguồn cùng pha là — phát biểu $d_2+d_1=k\lambda$ là đúng.
\end{itemchoice}
}
\end{ex}

% ===== Câu 42 =====
% ID: L11C2B12-11-TL
% Mức: VD
\dangbt{Nhận biết hiện tượng và điều kiện giao thoa sóng}
\begin{ex}
% ID: L11C2B12-11-TL
% Mức: VD
Trong dạng “Nhận biết hiện tượng và điều kiện giao thoa sóng”, Trình bày cách xác định bước sóng từ ảnh các đường cực đại trên đoạn nối hai nguồn.
\choiceTF

\loigiai{
Đo khoảng cách $i$ giữa hai cực đại liên tiếp trên đoạn nối nguồn; do $i=\lambda/2$ nên tính $\lambda=2i$.
}
\end{ex}

% ===== Câu 43 =====
% ID: L11C2B12-11-TLN
% Mức: VD
\begin{ex}
% ID: L11C2B12-11-TLN
% Mức: VD
Trong dạng “Nhận biết hiện tượng và điều kiện giao thoa sóng”, Hai nguồn cùng pha cách nhau $12\,\text{cm}$, $\lambda=3$ cm. Giá trị bậc cực đại lớn nhất trên đoạn nối nguồn là bao nhiêu?
\shortans{4}
\loigiai{
$|k|\le a/\lambda=12/3=4$, nên bậc lớn nhất là 4.
}
\end{ex}

% ===== Câu 44 =====
% ID: L11C2B12-11-TN
% Mức: NB
\dangbt{Chưa có dạng}
\begin{ex}
% ID: L11C2B12-11-TN
% Mức: NB
Hiện tượng giao thoa sóng là hiện tượng $\nguon{ly11 kntt.pdf}$
\choice
{giao nhau của hai sóng tại một điểm trong môi trường.}
{tổng hợp của hai dao động.}
{tạo thành các gợn lồi lõm.}
{\True hai sóng khi gặp nhau có những điểm cường độ sóng luôn tăng cường hoặc triệt tiêu nhau. 12.2. Hai nguồn kết hợp là hai nguồn có A. cùng biên độ. B. cùng tần số. C. cùng pha ban đầu. D. cùng tần số và hiệu số pha không đổi theo thời gian.}
\loigiai{
Chọn D.

Hiện tượng giao thoa sóng là hiện tượng hai sóng khi gặp nhau có những điểm cường độ sóng luôn tăng cường hoặc triệt tiêu nhau.
}
\end{ex}

% ===== Câu 45 =====
% ID: L11C2B12-12-DS
% Mức: TH
\dangbt{Tính hiệu đường đi và bậc giao thoa}
\begin{ex}
% ID: L11C2B12-12-DS
% Mức: TH
Xét các phát biểu sau trong dạng “Tính hiệu đường đi và bậc giao thoa”.
\choiceTF
{\True Điều kiện để tại $M$ có cực tiểu giao thoa của hai nguồn cùng pha là — phát biểu đúng là: $d_2-d_1=k\lambda$.}
{Điều kiện để tại $M$ có cực tiểu giao thoa của hai nguồn cùng pha là — phát biểu $d_2-d_1=(k+0,5)\lambda$ là đúng.}
{\True Hai nguồn cùng pha, $\lambda=2$ cm. Điểm $M$ có $d_2-d_1=6$ cm là — phát biểu đúng là: cực đại bậc 3.}
{Hai nguồn cùng pha, $\lambda=2$ cm. Điểm $M$ có $d_2-d_1=6$ cm là — phát biểu cực đại bậc 2 là đúng.}
\loigiai{
\begin{itemchoice}
\item (Đúng) Điều kiện để tại $M$ có cực tiểu giao thoa của hai nguồn cùng pha là — phát biểu đúng là: $d_2-d_1=k\lambda$. (Vì): Với hai nguồn cùng pha, cực tiểu khi hiệu đường đi bằng nửa nguyên lần bước sóng. b) Sai: phương án nêu không phù hợp. c) Đúng: Ta có $
\item (Sai) Điều kiện để tại $M$ có cực tiểu giao thoa của hai nguồn cùng pha là — phát biểu $d_2-d_1=(k+0,5)\lambda$ là đúng.
\item (Đúng) Hai nguồn cùng pha, $\lambda=2$ cm. Điểm $M$ có $d_2-d_1=6$ cm là — phát biểu đúng là: cực đại bậc 3.
\item (Sai) Hai nguồn cùng pha, $\lambda=2$ cm. Điểm $M$ có $d_2-d_1=6$ cm là — phát biểu cực đại bậc 2 là đúng. (Vì): _1)/\lambda=6/2=3$, nên $M$ thuộc cực đại bậc 3. d) Sai: phương án nêu không phù hợp
\end{itemchoice}
}
\end{ex}

% ===== Câu 46 =====
% ID: L11C2B12-12-TL
% Mức: VD
\dangbt{Nhận biết hiện tượng và điều kiện giao thoa sóng}
\begin{ex}
% ID: L11C2B12-12-TL
% Mức: VD
Trong dạng “Nhận biết hiện tượng và điều kiện giao thoa sóng”, Hai nguồn cùng pha cách nhau $16\,\text{cm}$, bước sóng $4\,\text{cm}$. Xác định số cực đại trên đoạn nối hai nguồn.
\choiceTF

\loigiai{
Điều kiện $|k|\le16/4=4$. Có các giá trị $k=-4,-3,\ldots,4$, tổng cộng 9 cực đại.
}
\end{ex}

% ===== Câu 47 =====
% ID: L11C2B12-12-TLN
% Mức: VDC
\begin{ex}
% ID: L11C2B12-12-TLN
% Mức: VDC
Trong dạng “Nhận biết hiện tượng và điều kiện giao thoa sóng”, Tại $M$, $d_1=18$ cm, $d_2=25$ cm. Tính độ lớn hiệu đường đi theo cm.
\shortans{7}
\loigiai{
$\Delta d=|d_2-d_1|=|25-18|=7$ cm.
}
\end{ex}

% ===== Câu 48 =====
% ID: L11C2B12-12-TN
% Mức: NB
\dangbt{Chưa có dạng}
\begin{ex}
% ID: L11C2B12-12-TN
% Mức: NB
Hai sóng phát ra từ hai nguồn kết hợp. Cực đại giao thoa nằm tại các điểm có hiệu khoảng cách tới hai nguồn sóng bằng $\nguon{ly11 kntt.pdf}$
\choice
{một ước số của bước sóng.}
{\True một bội số nguyên của bước sóng.}
{một bội số lẻ của nửa bước sóng.}
{một ước số của nửa bước sóng.}
\loigiai{
Chọn B.

Cực đại giao thoa nằm tại các điểm có hiệu khoảng cách tới hai nguồn sóng bằng một bội số nguyên của bước sóng.
}
\end{ex}

% ===== Câu 49 =====
% ID: L11C2B12-13-DS
% Mức: VD
\dangbt{Tính hiệu đường đi và bậc giao thoa}
\begin{ex}
% ID: L11C2B12-13-DS
% Mức: VD
Xét các phát biểu sau trong dạng “Tính hiệu đường đi và bậc giao thoa”.
\choiceTF
{\True Hai nguồn cùng pha, $\lambda=4$ cm. Điểm $N$ có $d_2-d_1=6$ cm là — phát biểu đúng là: cực đại.}
{Hai nguồn cùng pha, $\lambda=4$ cm. Điểm $N$ có $d_2-d_1=6$ cm là — phát biểu cực tiểu là đúng.}
{\True Trên đoạn nối hai nguồn cùng pha, khoảng cách giữa hai cực đại liên tiếp bằng — phát biểu đúng là: $\lambda/2$.}
{Trên đoạn nối hai nguồn cùng pha, khoảng cách giữa hai cực đại liên tiếp bằng — phát biểu $2\lambda$ là đúng.}
\loigiai{
\begin{itemchoice}
\item (Đúng) Hai nguồn cùng pha, $\lambda=4$ cm. Điểm $N$ có $d_2-d_1=6$ cm là — phát biểu đúng là: cực đại. (Vì): $6/4=1,5$, là nửa nguyên nên $N$ thuộc cực tiểu. b) Sai: phương án nêu không phù hợp. c) Đúng: Trên đoạn nối hai nguồn, hiệu đường đi đổi hai lần theo li độ nên hai cực đại liên tiếp cách nhau $\lambda/2$. d) Sai: phương án nêu không phù hợp
\item (Sai) Hai nguồn cùng pha, $\lambda=4$ cm. Điểm $N$ có $d_2-d_1=6$ cm là — phát biểu cực tiểu là đúng.
\item (Đúng) Trên đoạn nối hai nguồn cùng pha, khoảng cách giữa hai cực đại liên tiếp bằng — phát biểu đúng là: $\lambda/2$.
\item (Sai) Trên đoạn nối hai nguồn cùng pha, khoảng cách giữa hai cực đại liên tiếp bằng — phát biểu $2\lambda$ là đúng.
\end{itemchoice}
}
\end{ex}

% ===== Câu 50 =====
% ID: L11C2B12-13-TL
% Mức: VDC
\dangbt{Nhận biết hiện tượng và điều kiện giao thoa sóng}
\begin{ex}
% ID: L11C2B12-13-TL
% Mức: VDC
Trong dạng “Nhận biết hiện tượng và điều kiện giao thoa sóng”, Phân biệt cực đại và cực tiểu giao thoa của hai nguồn cùng pha bằng hiệu đường đi.
\choiceTF

\loigiai{
Cực đại: $\Delta d=k\lambda$. Cực tiểu: $\Delta d=(k+1/2)\lambda$, với $k$ nguyên.
}
\end{ex}

% ===== Câu 51 =====
% ID: L11C2B12-13-TLN
% Mức: TH
\dangbt{Tính hiệu đường đi và bậc giao thoa}
\begin{ex}
% ID: L11C2B12-13-TLN
% Mức: TH
Trong dạng “Tính hiệu đường đi và bậc giao thoa”, Hai nguồn cùng pha, $\lambda=3$ cm. Tại $M$, hiệu đường đi là $12\,\text{cm}$. Xác định bậc cực đại.
\shortans{4}
\loigiai{
$k=(d_2-d_1)/\lambda=12/3=4$.
}
\end{ex}

% ===== Câu 52 =====
% ID: L11C2B12-13-TN
% Mức: NB
\dangbt{Chưa có dạng}
\begin{ex}
% ID: L11C2B12-13-TN
% Mức: NB
Một trong hai khe của thí nghiệm Young được làm mờ sao cho nó chỉ truyền ánh sáng được bằng 1/2 cường độ sáng của khe còn lại. Kết quả là $\nguon{ly11 kntt.pdf}$
\choice
{vân giao thoa biến mất.}
{vân giao thoa tối đi.}
{vạch sáng trở nên sáng hơn và vạch tối thì tối hơn.}
{\True vạch tối sáng hơn và vạch sáng tối hơn.}
\loigiai{
Chọn D.

Cường độ sáng giảm thì vạch tối sáng hơn và vạch sáng tối hơn.
}
\end{ex}

% ===== Câu 53 =====
% ID: L11C2B12-14-DS
% Mức: VD
\dangbt{Tính hiệu đường đi và bậc giao thoa}
\begin{ex}
% ID: L11C2B12-14-DS
% Mức: VD
Xét các phát biểu sau trong dạng “Tính hiệu đường đi và bậc giao thoa”.
\choiceTF
{\True Hiện tượng giao thoa là bằng chứng rõ ràng của — phát biểu đúng là: tính chất sóng.}
{Hiện tượng giao thoa là bằng chứng rõ ràng của — phát biểu tính chất hạt là đúng.}
{\True Hai nguồn cùng pha cách nhau $10\,\text{cm}$, $\lambda=2$ cm. Số cực đại trên đoạn nối hai nguồn, kể cả hai nguồn nếu thỏa điều kiện, là — phát biểu đúng là: 9.}
{Hai nguồn cùng pha cách nhau $10\,\text{cm}$, $\lambda=2$ cm. Số cực đại trên đoạn nối hai nguồn, kể cả hai nguồn nếu thỏa điều kiện, là — phát biểu 11 là đúng.}
\loigiai{
\begin{itemchoice}
\item (Đúng) Hiện tượng giao thoa là bằng chứng rõ ràng của — phát biểu đúng là: tính chất sóng. (Vì): Giao thoa là hiện tượng đặc trưng của sóng. b) Sai: phương án nêu không phù hợp. c) Đúng: Các bậc nguyên thỏa $|k|\le 5$, nên có 11 giá trị từ -5 đến 5. d) Sai: phương án nêu không phù hợp
\item (Sai) Hiện tượng giao thoa là bằng chứng rõ ràng của — phát biểu tính chất hạt là đúng.
\item (Đúng) Hai nguồn cùng pha cách nhau $10\,\text{cm}$, $\lambda=2$ cm. Số cực đại trên đoạn nối hai nguồn, kể cả hai nguồn nếu thỏa điều kiện, là — phát biểu đúng là: 9.
\item (Sai) Hai nguồn cùng pha cách nhau $10\,\text{cm}$, $\lambda=2$ cm. Số cực đại trên đoạn nối hai nguồn, kể cả hai nguồn nếu thỏa điều kiện, là — phát biểu 11 là đúng.
\end{itemchoice}
}
\end{ex}

% ===== Câu 54 =====
% ID: L11C2B12-14-TL
% Mức: TH
\begin{ex}
% ID: L11C2B12-14-TL
% Mức: TH
Trong dạng “Tính hiệu đường đi và bậc giao thoa”, Nêu điều kiện để hai nguồn sóng là hai nguồn kết hợp.
\choiceTF

\loigiai{
Hai nguồn dao động cùng phương, cùng tần số và có độ lệch pha không đổi theo thời gian.
}
\end{ex}

% ===== Câu 55 =====
% ID: L11C2B12-14-TLN
% Mức: TH
\begin{ex}
% ID: L11C2B12-14-TLN
% Mức: TH
Trong dạng “Tính hiệu đường đi và bậc giao thoa”, Hai nguồn cùng pha, $\lambda=4$ cm. Hiệu đường đi tại cực tiểu thứ hai là bao nhiêu theo cm?
\shortans{6}
\loigiai{
Cực tiểu thứ hai ứng với $|d_2-d_1|=(1+0,5)\lambda=1,5\cdot4=6$ cm.
}
\end{ex}

% ===== Câu 56 =====
% ID: L11C2B12-14-TN
% Mức: NB
\dangbt{Chưa có dạng}
\begin{ex}
% ID: L11C2B12-14-TN
% Mức: NB
Trong thí nghiệm Young về giao thoa với ánh sáng đơn sắc, khoảng cách giữa hai khe là $0,15 \mathrm{mm}$, khoảng cách giữa mặt phẳng chứa hai khe và màn quan sát là $2\,\text{m}$. Ânh sáng sử dụng trong thi nghiệm là ánh sáng đơn sắc màu vàng có bước sóng 0,58 $\mu$ m. Vị trí vân sáng bậc 3 trên màn quan sát cách vân trung tâm một khoảng là $\nguon{ly11 kntt.pdf}$
\choice
{$0,232 \cdot  10^{-3} m.$}
{$0,812 \cdot  10^{-3} m.$}
{\True $2,23 \cdot  10^{-3} m.$}
{$8,12 \cdot  10^{-3} m.$}
\loigiai{
Chọn C.

x = ki = k $\dfrac{\lambda D}{a} = 3\dfrac{0,58 \cdot 10^{-6} \cdot 2}{0,15 \cdot 10^{-3}}=0,232\,\mathrm{mm}$.
}
\end{ex}

% ===== Câu 57 =====
% ID: L11C2B12-15-DS
% Mức: VDC
\dangbt{Tính hiệu đường đi và bậc giao thoa}
\begin{ex}
% ID: L11C2B12-15-DS
% Mức: VDC
Xét các phát biểu sau trong dạng “Tính hiệu đường đi và bậc giao thoa”.
\choiceTF
{\True Tại điểm có cực đại giao thoa, hai sóng thành phần đến đó — phát biểu đúng là: cùng pha.}
{Tại điểm có cực đại giao thoa, hai sóng thành phần đến đó — phát biểu ngược pha là đúng.}
{\True Nếu tần số nguồn tăng gấp đôi còn tốc độ truyền sóng không đổi thì bước sóng — phát biểu đúng là: giảm một nửa.}
{Nếu tần số nguồn tăng gấp đôi còn tốc độ truyền sóng không đổi thì bước sóng — phát biểu không đổi là đúng.}
\loigiai{
\begin{itemchoice}
\item (Đúng) Tại điểm có cực đại giao thoa, hai sóng thành phần đến đó — phát biểu đúng là: cùng pha. (Vì): Cực đại xuất hiện khi hai dao động thành phần tăng cường nhau, tức cùng pha. b) Sai: phương án nêu không phù hợp. c) Đúng: Từ $\lambda=v/f$, khi $f$ tăng gấp đôi thì $\lambda$ giảm một nửa. d) Sai: phương án nêu không phù hợp
\item (Sai) Tại điểm có cực đại giao thoa, hai sóng thành phần đến đó — phát biểu ngược pha là đúng.
\item (Đúng) Nếu tần số nguồn tăng gấp đôi còn tốc độ truyền sóng không đổi thì bước sóng — phát biểu đúng là: giảm một nửa.
\item (Sai) Nếu tần số nguồn tăng gấp đôi còn tốc độ truyền sóng không đổi thì bước sóng — phát biểu không đổi là đúng.
\end{itemchoice}
}
\end{ex}

% ===== Câu 58 =====
% ID: L11C2B12-15-TL
% Mức: TH
\begin{ex}
% ID: L11C2B12-15-TL
% Mức: TH
Trong dạng “Tính hiệu đường đi và bậc giao thoa”, Giải thích vì sao ở cực đại giao thoa biên độ dao động lớn.
\choiceTF

\loigiai{
Tại cực đại, hai sóng đến cùng pha nên li độ tổng hợp được cộng đại số và biên độ đạt giá trị lớn.
}
\end{ex}

% ===== Câu 59 =====
% ID: L11C2B12-15-TLN
% Mức: VD
\begin{ex}
% ID: L11C2B12-15-TLN
% Mức: VD
Trong dạng “Tính hiệu đường đi và bậc giao thoa”, Trên đoạn nối hai nguồn, hai cực đại liên tiếp cách nhau $2,5\,\text{cm}$. Tính bước sóng theo cm.
\shortans{5}
\loigiai{
Khoảng cách hai cực đại liên tiếp trên đoạn nối nguồn bằng $\lambda/2$, nên $\lambda=5$ cm.
}
\end{ex}

% ===== Câu 60 =====
% ID: L11C2B12-15-TN
% Mức: NB
\dangbt{Chưa có dạng}
\begin{ex}
% ID: L11C2B12-15-TN
% Mức: NB
Trong thí nghiệm Young về giao thoa ánh sáng, giữa hai điểm $P$ và $Q$ trên màn cách nhau $9 \mathrm{mm}$ chỉ có 5 vân sáng mà tại $P$ là một trong 5 vân sáng đó, còn tại $Q$ là vị trí của vân tối. Vị trí vân tối thứ 2 cách vân trung tâm một khoảng là $\nguon{ly11 kntt.pdf}$
\choice
{$0,5 \cdot  10^{-3} m.$}
{$5 \cdot  10^{-3} m.$}
{\True $3 \cdot  10^{-3} m.$}
{$0,3 \cdot  10^{-3} m.$}
\loigiai{
Chọn C.

Hai điểm $P$ và $Q$ trên màn cách nhau $9\,\mathrm{mm}$ chỉ có 5 vân sáng mà tại $P$ là một trong 5 vân sáng đó, còn tại $Q$ là vị trí của vân tối nên PQ = 4,5i $\Rightarrow$ i = $2\,\mathrm{mm}$

Vị trí vân tối thứ 2: x_(t2)= 1,5i = $3\,\mathrm{mm}$
}
\end{ex}

% ===== Câu 61 =====
% ID: L11C2B12-16-DS
% Mức: TH
\dangbt{Vận dụng giao thoa sóng trong thực tế}
\begin{ex}
% ID: L11C2B12-16-DS
% Mức: TH
Xét các phát biểu sau trong dạng “Vận dụng giao thoa sóng trong thực tế”.
\choiceTF
{\True Hai nguồn kết hợp là hai nguồn dao động — phát biểu đúng là: cùng phương, cùng tần số và có độ lệch pha không đổi.}
{Hai nguồn kết hợp là hai nguồn dao động — phát biểu cùng biên độ nhưng khác tần số là đúng.}
{\True Điều kiện để tại $M$ có cực đại giao thoa của hai nguồn cùng pha là — phát biểu đúng là: $d_2-d_1=k\lambda$.}
{Điều kiện để tại $M$ có cực đại giao thoa của hai nguồn cùng pha là — phát biểu $d_2+d_1=k\lambda$ là đúng.}
\loigiai{
\begin{itemchoice}
\item (Đúng) Hai nguồn kết hợp là hai nguồn dao động — phát biểu đúng là: cùng phương, cùng tần số và có độ lệch pha không đổi. (Vì): Hai nguồn kết hợp phải cùng phương, cùng tần số và có độ lệch pha không đổi. b) Sai: phương án nêu không phù hợp. c) Đúng: Với hai nguồn cùng pha, cực đại khi hiệu đường đi bằng số nguyên lần bước sóng. d) Sai: phương án nêu không phù hợp
\item (Sai) Hai nguồn kết hợp là hai nguồn dao động — phát biểu cùng biên độ nhưng khác tần số là đúng.
\item (Đúng) Điều kiện để tại $M$ có cực đại giao thoa của hai nguồn cùng pha là — phát biểu đúng là: $d_2-d_1=k\lambda$.
\item (Sai) Điều kiện để tại $M$ có cực đại giao thoa của hai nguồn cùng pha là — phát biểu $d_2+d_1=k\lambda$ là đúng.
\end{itemchoice}
}
\end{ex}

% ===== Câu 62 =====
% ID: L11C2B12-16-TL
% Mức: VD
\dangbt{Tính hiệu đường đi và bậc giao thoa}
\begin{ex}
% ID: L11C2B12-16-TL
% Mức: VD
Trong dạng “Tính hiệu đường đi và bậc giao thoa”, Hai nguồn cùng pha có $\lambda=2$ cm. Xét điểm $M$ với $d_1=10$ cm, $d_2=15$ cm. Xác định trạng thái giao thoa tại M.
\choiceTF

\loigiai{
Hiệu đường đi $\Delta d=5\,\mathrm{cm}=2,5\lambda$, nên hai sóng ngược pha và $M$ thuộc cực tiểu.
}
\end{ex}

% ===== Câu 63 =====
% ID: L11C2B12-16-TLN
% Mức: VD
\begin{ex}
% ID: L11C2B12-16-TLN
% Mức: VD
Trong dạng “Tính hiệu đường đi và bậc giao thoa”, Sóng có tốc độ $40\,\text{cm}$ /s và tần số $10\,\text{Hz}$. Tính bước sóng theo cm.
\shortans{4}
\loigiai{
$\lambda=v/f=40/10=4$ cm.
}
\end{ex}

% ===== Câu 64 =====
% ID: L11C2B12-16-TN
% Mức: NB
\dangbt{Chưa có dạng}
\begin{ex}
% ID: L11C2B12-16-TN
% Mức: NB
Trong thí nghiệm Young về giao thoa với ánh sáng đơn sắc, khoảng cách giữa hai khe là $0,15 \mathrm{mm}$, khoảng cách giữa mặt phẳng chứa hai khe và màn quan sát là $1,5\,\text{m}$. Khoảng cách giữa 5 vân sáng liên tiếp là $36 \mathrm{mm}$. Bước sóng của ánh sáng dùng trong thí nghiệm này là $\nguon{ly11 kntt.pdf}$
\choice
{0,60 $\mu$ m.}
{0,40 $\mu$ m.}
{\True 0,48 $\mu$ m.}
{0,76 $\mu$ m.}
\loigiai{
Chọn Không có đáp án phù hợp.

Đáp án đúng là không có

$\left. i = \dfrac{\text{\Delta }s}{n - 1} = \dfrac{36}{5 - 1} = 9\text{mm}\Rightarrow\lambda = \dfrac{ai}{D} = \dfrac{0,15\cdot 10^{-3} \cdot 9 \cdot 10^{-3}}{1,5} = 9 \cdot 10^{-7}\text{m} \right.$.
}
\end{ex}

% ===== Câu 65 =====
% ID: L11C2B12-17-DS
% Mức: TH
\dangbt{Vận dụng giao thoa sóng trong thực tế}
\begin{ex}
% ID: L11C2B12-17-DS
% Mức: TH
Xét các phát biểu sau trong dạng “Vận dụng giao thoa sóng trong thực tế”.
\choiceTF
{\True Điều kiện để tại $M$ có cực tiểu giao thoa của hai nguồn cùng pha là — phát biểu đúng là: $d_2-d_1=k\lambda$.}
{Điều kiện để tại $M$ có cực tiểu giao thoa của hai nguồn cùng pha là — phát biểu $d_2-d_1=(k+0,5)\lambda$ là đúng.}
{\True Hai nguồn cùng pha, $\lambda=2$ cm. Điểm $M$ có $d_2-d_1=6$ cm là — phát biểu đúng là: cực đại bậc 3.}
{Hai nguồn cùng pha, $\lambda=2$ cm. Điểm $M$ có $d_2-d_1=6$ cm là — phát biểu cực đại bậc 2 là đúng.}
\loigiai{
\begin{itemchoice}
\item (Đúng) Điều kiện để tại $M$ có cực tiểu giao thoa của hai nguồn cùng pha là — phát biểu đúng là: $d_2-d_1=k\lambda$. (Vì): Với hai nguồn cùng pha, cực tiểu khi hiệu đường đi bằng nửa nguyên lần bước sóng. b) Sai: phương án nêu không phù hợp. c) Đúng: Ta có $
\item (Sai) Điều kiện để tại $M$ có cực tiểu giao thoa của hai nguồn cùng pha là — phát biểu $d_2-d_1=(k+0,5)\lambda$ là đúng.
\item (Đúng) Hai nguồn cùng pha, $\lambda=2$ cm. Điểm $M$ có $d_2-d_1=6$ cm là — phát biểu đúng là: cực đại bậc 3.
\item (Sai) Hai nguồn cùng pha, $\lambda=2$ cm. Điểm $M$ có $d_2-d_1=6$ cm là — phát biểu cực đại bậc 2 là đúng. (Vì): _1)/\lambda=6/2=3$, nên $M$ thuộc cực đại bậc 3. d) Sai: phương án nêu không phù hợp
\end{itemchoice}
}
\end{ex}

% ===== Câu 66 =====
% ID: L11C2B12-17-TL
% Mức: VD
\dangbt{Tính hiệu đường đi và bậc giao thoa}
\begin{ex}
% ID: L11C2B12-17-TL
% Mức: VD
Trong dạng “Tính hiệu đường đi và bậc giao thoa”, Trình bày cách xác định bước sóng từ ảnh các đường cực đại trên đoạn nối hai nguồn.
\choiceTF

\loigiai{
Đo khoảng cách $i$ giữa hai cực đại liên tiếp trên đoạn nối nguồn; do $i=\lambda/2$ nên tính $\lambda=2i$.
}
\end{ex}

% ===== Câu 67 =====
% ID: L11C2B12-17-TLN
% Mức: VD
\begin{ex}
% ID: L11C2B12-17-TLN
% Mức: VD
Trong dạng “Tính hiệu đường đi và bậc giao thoa”, Hai nguồn cùng pha cách nhau $12\,\text{cm}$, $\lambda=3$ cm. Giá trị bậc cực đại lớn nhất trên đoạn nối nguồn là bao nhiêu?
\shortans{4}
\loigiai{
$|k|\le a/\lambda=12/3=4$, nên bậc lớn nhất là 4.
}
\end{ex}

% ===== Câu 68 =====
% ID: L11C2B12-17-TN
% Mức: NB
\dangbt{Nhận biết hiện tượng và điều kiện giao thoa sóng}
\begin{ex}
% ID: L11C2B12-17-TN
% Mức: NB
Trong dạng “Nhận biết hiện tượng và điều kiện giao thoa sóng”, Hai nguồn kết hợp là hai nguồn dao động
\choice
{\True cùng phương, cùng tần số và có độ lệch pha không đổi}
{cùng biên độ nhưng khác tần số}
{cùng pha ở mọi thời điểm dù khác tần số}
{vuông pha và khác phương}
\loigiai{
Đáp án A. Hai nguồn kết hợp phải cùng phương, cùng tần số và có độ lệch pha không đổi.
}
\end{ex}

% ===== Câu 69 =====
% ID: L11C2B12-18-DS
% Mức: VD
\dangbt{Vận dụng giao thoa sóng trong thực tế}
\begin{ex}
% ID: L11C2B12-18-DS
% Mức: VD
Xét các phát biểu sau trong dạng “Vận dụng giao thoa sóng trong thực tế”.
\choiceTF
{\True Hai nguồn cùng pha, $\lambda=4$ cm. Điểm $N$ có $d_2-d_1=6$ cm là — phát biểu đúng là: cực đại.}
{Hai nguồn cùng pha, $\lambda=4$ cm. Điểm $N$ có $d_2-d_1=6$ cm là — phát biểu cực tiểu là đúng.}
{\True Trên đoạn nối hai nguồn cùng pha, khoảng cách giữa hai cực đại liên tiếp bằng — phát biểu đúng là: $\lambda/2$.}
{Trên đoạn nối hai nguồn cùng pha, khoảng cách giữa hai cực đại liên tiếp bằng — phát biểu $2\lambda$ là đúng.}
\loigiai{
\begin{itemchoice}
\item (Đúng) Hai nguồn cùng pha, $\lambda=4$ cm. Điểm $N$ có $d_2-d_1=6$ cm là — phát biểu đúng là: cực đại. (Vì): $6/4=1,5$, là nửa nguyên nên $N$ thuộc cực tiểu. b) Sai: phương án nêu không phù hợp. c) Đúng: Trên đoạn nối hai nguồn, hiệu đường đi đổi hai lần theo li độ nên hai cực đại liên tiếp cách nhau $\lambda/2$. d) Sai: phương án nêu không phù hợp
\item (Sai) Hai nguồn cùng pha, $\lambda=4$ cm. Điểm $N$ có $d_2-d_1=6$ cm là — phát biểu cực tiểu là đúng.
\item (Đúng) Trên đoạn nối hai nguồn cùng pha, khoảng cách giữa hai cực đại liên tiếp bằng — phát biểu đúng là: $\lambda/2$.
\item (Sai) Trên đoạn nối hai nguồn cùng pha, khoảng cách giữa hai cực đại liên tiếp bằng — phát biểu $2\lambda$ là đúng.
\end{itemchoice}
}
\end{ex}

% ===== Câu 70 =====
% ID: L11C2B12-18-TL
% Mức: VD
\dangbt{Tính hiệu đường đi và bậc giao thoa}
\begin{ex}
% ID: L11C2B12-18-TL
% Mức: VD
Trong dạng “Tính hiệu đường đi và bậc giao thoa”, Hai nguồn cùng pha cách nhau $16\,\text{cm}$, bước sóng $4\,\text{cm}$. Xác định số cực đại trên đoạn nối hai nguồn.
\choiceTF

\loigiai{
Điều kiện $|k|\le16/4=4$. Có các giá trị $k=-4,-3,\ldots,4$, tổng cộng 9 cực đại.
}
\end{ex}

% ===== Câu 71 =====
% ID: L11C2B12-18-TLN
% Mức: VDC
\begin{ex}
% ID: L11C2B12-18-TLN
% Mức: VDC
Trong dạng “Tính hiệu đường đi và bậc giao thoa”, Tại $M$, $d_1=18$ cm, $d_2=25$ cm. Tính độ lớn hiệu đường đi theo cm.
\shortans{7}
\loigiai{
$\Delta d=|d_2-d_1|=|25-18|=7$ cm.
}
\end{ex}

% ===== Câu 72 =====
% ID: L11C2B12-18-TN
% Mức: NB
\dangbt{Nhận biết hiện tượng và điều kiện giao thoa sóng}
\begin{ex}
% ID: L11C2B12-18-TN
% Mức: NB
Trong dạng “Nhận biết hiện tượng và điều kiện giao thoa sóng”, Điều kiện để tại $M$ có cực đại giao thoa của hai nguồn cùng pha là
\choice
{\True $d_2-d_1=k\lambda$}
{$d_2-d_1=(k+0,5)\lambda$}
{$d_2+d_1=k\lambda$}
{$d_2-d_1=\lambda/4$}
\loigiai{
Đáp án A. Với hai nguồn cùng pha, cực đại khi hiệu đường đi bằng số nguyên lần bước sóng.
}
\end{ex}

% ===== Câu 73 =====
% ID: L11C2B12-19-DS
% Mức: VD
\dangbt{Vận dụng giao thoa sóng trong thực tế}
\begin{ex}
% ID: L11C2B12-19-DS
% Mức: VD
Xét các phát biểu sau trong dạng “Vận dụng giao thoa sóng trong thực tế”.
\choiceTF
{\True Hiện tượng giao thoa là bằng chứng rõ ràng của — phát biểu đúng là: tính chất sóng.}
{Hiện tượng giao thoa là bằng chứng rõ ràng của — phát biểu tính chất hạt là đúng.}
{\True Hai nguồn cùng pha cách nhau $10\,\text{cm}$, $\lambda=2$ cm. Số cực đại trên đoạn nối hai nguồn, kể cả hai nguồn nếu thỏa điều kiện, là — phát biểu đúng là: 9.}
{Hai nguồn cùng pha cách nhau $10\,\text{cm}$, $\lambda=2$ cm. Số cực đại trên đoạn nối hai nguồn, kể cả hai nguồn nếu thỏa điều kiện, là — phát biểu 11 là đúng.}
\loigiai{
\begin{itemchoice}
\item (Đúng) Hiện tượng giao thoa là bằng chứng rõ ràng của — phát biểu đúng là: tính chất sóng. (Vì): Giao thoa là hiện tượng đặc trưng của sóng. b) Sai: phương án nêu không phù hợp. c) Đúng: Các bậc nguyên thỏa $|k|\le 5$, nên có 11 giá trị từ -5 đến 5. d) Sai: phương án nêu không phù hợp
\item (Sai) Hiện tượng giao thoa là bằng chứng rõ ràng của — phát biểu tính chất hạt là đúng.
\item (Đúng) Hai nguồn cùng pha cách nhau $10\,\text{cm}$, $\lambda=2$ cm. Số cực đại trên đoạn nối hai nguồn, kể cả hai nguồn nếu thỏa điều kiện, là — phát biểu đúng là: 9.
\item (Sai) Hai nguồn cùng pha cách nhau $10\,\text{cm}$, $\lambda=2$ cm. Số cực đại trên đoạn nối hai nguồn, kể cả hai nguồn nếu thỏa điều kiện, là — phát biểu 11 là đúng.
\end{itemchoice}
}
\end{ex}

% ===== Câu 74 =====
% ID: L11C2B12-19-TL
% Mức: VDC
\dangbt{Tính hiệu đường đi và bậc giao thoa}
\begin{ex}
% ID: L11C2B12-19-TL
% Mức: VDC
Trong dạng “Tính hiệu đường đi và bậc giao thoa”, Phân biệt cực đại và cực tiểu giao thoa của hai nguồn cùng pha bằng hiệu đường đi.
\choiceTF

\loigiai{
Cực đại: $\Delta d=k\lambda$. Cực tiểu: $\Delta d=(k+1/2)\lambda$, với $k$ nguyên.
}
\end{ex}

% ===== Câu 75 =====
% ID: L11C2B12-19-TLN
% Mức: TH
\dangbt{Vận dụng giao thoa sóng trong thực tế}
\begin{ex}
% ID: L11C2B12-19-TLN
% Mức: TH
Trong dạng “Vận dụng giao thoa sóng trong thực tế”, Hai nguồn cùng pha, $\lambda=3$ cm. Tại $M$, hiệu đường đi là $12\,\text{cm}$. Xác định bậc cực đại.
\shortans{4}
\loigiai{
$k=(d_2-d_1)/\lambda=12/3=4$.
}
\end{ex}

% ===== Câu 76 =====
% ID: L11C2B12-19-TN
% Mức: TH
\dangbt{Nhận biết hiện tượng và điều kiện giao thoa sóng}
\begin{ex}
% ID: L11C2B12-19-TN
% Mức: TH
Trong dạng “Nhận biết hiện tượng và điều kiện giao thoa sóng”, Điều kiện để tại $M$ có cực tiểu giao thoa của hai nguồn cùng pha là
\choice
{$d_2-d_1=k\lambda$}
{\True $d_2-d_1=(k+0,5)\lambda$}
{$d_2+d_1=(k+0,5)\lambda$}
{$d_2-d_1=2k\lambda$}
\loigiai{
Đáp án B. Với hai nguồn cùng pha, cực tiểu khi hiệu đường đi bằng nửa nguyên lần bước sóng.
}
\end{ex}

% ===== Câu 77 =====
% ID: L11C2B12-20-DS
% Mức: VDC
\dangbt{Vận dụng giao thoa sóng trong thực tế}
\begin{ex}
% ID: L11C2B12-20-DS
% Mức: VDC
Xét các phát biểu sau trong dạng “Vận dụng giao thoa sóng trong thực tế”.
\choiceTF
{\True Tại điểm có cực đại giao thoa, hai sóng thành phần đến đó — phát biểu đúng là: cùng pha.}
{Tại điểm có cực đại giao thoa, hai sóng thành phần đến đó — phát biểu ngược pha là đúng.}
{\True Nếu tần số nguồn tăng gấp đôi còn tốc độ truyền sóng không đổi thì bước sóng — phát biểu đúng là: giảm một nửa.}
{Nếu tần số nguồn tăng gấp đôi còn tốc độ truyền sóng không đổi thì bước sóng — phát biểu không đổi là đúng.}
\loigiai{
\begin{itemchoice}
\item (Đúng) Tại điểm có cực đại giao thoa, hai sóng thành phần đến đó — phát biểu đúng là: cùng pha. (Vì): Cực đại xuất hiện khi hai dao động thành phần tăng cường nhau, tức cùng pha. b) Sai: phương án nêu không phù hợp. c) Đúng: Từ $\lambda=v/f$, khi $f$ tăng gấp đôi thì $\lambda$ giảm một nửa. d) Sai: phương án nêu không phù hợp
\item (Sai) Tại điểm có cực đại giao thoa, hai sóng thành phần đến đó — phát biểu ngược pha là đúng.
\item (Đúng) Nếu tần số nguồn tăng gấp đôi còn tốc độ truyền sóng không đổi thì bước sóng — phát biểu đúng là: giảm một nửa.
\item (Sai) Nếu tần số nguồn tăng gấp đôi còn tốc độ truyền sóng không đổi thì bước sóng — phát biểu không đổi là đúng.
\end{itemchoice}
}
\end{ex}

% ===== Câu 78 =====
% ID: L11C2B12-20-TL
% Mức: TH
\begin{ex}
% ID: L11C2B12-20-TL
% Mức: TH
Trong dạng “Vận dụng giao thoa sóng trong thực tế”, Nêu điều kiện để hai nguồn sóng là hai nguồn kết hợp.
\choiceTF

\loigiai{
Hai nguồn dao động cùng phương, cùng tần số và có độ lệch pha không đổi theo thời gian.
}
\end{ex}

% ===== Câu 79 =====
% ID: L11C2B12-20-TLN
% Mức: TH
\begin{ex}
% ID: L11C2B12-20-TLN
% Mức: TH
Trong dạng “Vận dụng giao thoa sóng trong thực tế”, Hai nguồn cùng pha, $\lambda=4$ cm. Hiệu đường đi tại cực tiểu thứ hai là bao nhiêu theo cm?
\shortans{6}
\loigiai{
Cực tiểu thứ hai ứng với $|d_2-d_1|=(1+0,5)\lambda=1,5\cdot4=6$ cm.
}
\end{ex}

% ===== Câu 80 =====
% ID: L11C2B12-20-TN
% Mức: TH
\dangbt{Nhận biết hiện tượng và điều kiện giao thoa sóng}
\begin{ex}
% ID: L11C2B12-20-TN
% Mức: TH
Trong dạng “Nhận biết hiện tượng và điều kiện giao thoa sóng”, Hai nguồn cùng pha, $\lambda=2$ cm. Điểm $M$ có $d_2-d_1=6$ cm là
\choice
{\True cực đại bậc 3}
{cực tiểu thứ 3}
{cực đại bậc 2}
{không giao thoa}
\loigiai{
Đáp án A. Ta có $(d_2-d_1)/\lambda=6/2=3$, nên $M$ thuộc cực đại bậc 3.
}
\end{ex}

% ===== Câu 81 =====
% ID: L11C2B12-21-DS
% Mức: TH
\dangbt{Xác định cực đại, cực tiểu giao thoa}
\begin{ex}
% ID: L11C2B12-21-DS
% Mức: TH
Xét các phát biểu sau trong dạng “Xác định cực đại, cực tiểu giao thoa”.
\choiceTF
{\True Hai nguồn kết hợp là hai nguồn dao động — phát biểu đúng là: cùng phương, cùng tần số và có độ lệch pha không đổi.}
{Hai nguồn kết hợp là hai nguồn dao động — phát biểu cùng biên độ nhưng khác tần số là đúng.}
{\True Điều kiện để tại $M$ có cực đại giao thoa của hai nguồn cùng pha là — phát biểu đúng là: $d_2-d_1=k\lambda$.}
{Điều kiện để tại $M$ có cực đại giao thoa của hai nguồn cùng pha là — phát biểu $d_2+d_1=k\lambda$ là đúng.}
\loigiai{
\begin{itemchoice}
\item (Đúng) Hai nguồn kết hợp là hai nguồn dao động — phát biểu đúng là: cùng phương, cùng tần số và có độ lệch pha không đổi. (Vì): Hai nguồn kết hợp phải cùng phương, cùng tần số và có độ lệch pha không đổi. b) Sai: phương án nêu không phù hợp. c) Đúng: Với hai nguồn cùng pha, cực đại khi hiệu đường đi bằng số nguyên lần bước sóng. d) Sai: phương án nêu không phù hợp
\item (Sai) Hai nguồn kết hợp là hai nguồn dao động — phát biểu cùng biên độ nhưng khác tần số là đúng.
\item (Đúng) Điều kiện để tại $M$ có cực đại giao thoa của hai nguồn cùng pha là — phát biểu đúng là: $d_2-d_1=k\lambda$.
\item (Sai) Điều kiện để tại $M$ có cực đại giao thoa của hai nguồn cùng pha là — phát biểu $d_2+d_1=k\lambda$ là đúng.
\end{itemchoice}
}
\end{ex}

% ===== Câu 82 =====
% ID: L11C2B12-21-TL
% Mức: TH
\dangbt{Vận dụng giao thoa sóng trong thực tế}
\begin{ex}
% ID: L11C2B12-21-TL
% Mức: TH
Trong dạng “Vận dụng giao thoa sóng trong thực tế”, Giải thích vì sao ở cực đại giao thoa biên độ dao động lớn.
\choiceTF

\loigiai{
Tại cực đại, hai sóng đến cùng pha nên li độ tổng hợp được cộng đại số và biên độ đạt giá trị lớn.
}
\end{ex}

% ===== Câu 83 =====
% ID: L11C2B12-21-TLN
% Mức: VD
\begin{ex}
% ID: L11C2B12-21-TLN
% Mức: VD
Trong dạng “Vận dụng giao thoa sóng trong thực tế”, Trên đoạn nối hai nguồn, hai cực đại liên tiếp cách nhau $2,5\,\text{cm}$. Tính bước sóng theo cm.
\shortans{5}
\loigiai{
Khoảng cách hai cực đại liên tiếp trên đoạn nối nguồn bằng $\lambda/2$, nên $\lambda=5$ cm.
}
\end{ex}

% ===== Câu 84 =====
% ID: L11C2B12-21-TN
% Mức: TH
\dangbt{Nhận biết hiện tượng và điều kiện giao thoa sóng}
\begin{ex}
% ID: L11C2B12-21-TN
% Mức: TH
Trong dạng “Nhận biết hiện tượng và điều kiện giao thoa sóng”, Hai nguồn cùng pha, $\lambda=4$ cm. Điểm $N$ có $d_2-d_1=6$ cm là
\choice
{cực đại}
{\True cực tiểu}
{điểm đứng yên do không có sóng tới}
{không xác định}
\loigiai{
Đáp án B. $6/4=1,5$, là nửa nguyên nên $N$ thuộc cực tiểu.
}
\end{ex}

% ===== Câu 85 =====
% ID: L11C2B12-22-DS
% Mức: TH
\dangbt{Xác định cực đại, cực tiểu giao thoa}
\begin{ex}
% ID: L11C2B12-22-DS
% Mức: TH
Xét các phát biểu sau trong dạng “Xác định cực đại, cực tiểu giao thoa”.
\choiceTF
{\True Điều kiện để tại $M$ có cực tiểu giao thoa của hai nguồn cùng pha là — phát biểu đúng là: $d_2-d_1=k\lambda$.}
{Điều kiện để tại $M$ có cực tiểu giao thoa của hai nguồn cùng pha là — phát biểu $d_2-d_1=(k+0,5)\lambda$ là đúng.}
{\True Hai nguồn cùng pha, $\lambda=2$ cm. Điểm $M$ có $d_2-d_1=6$ cm là — phát biểu đúng là: cực đại bậc 3.}
{Hai nguồn cùng pha, $\lambda=2$ cm. Điểm $M$ có $d_2-d_1=6$ cm là — phát biểu cực đại bậc 2 là đúng.}
\loigiai{
\begin{itemchoice}
\item (Đúng) Điều kiện để tại $M$ có cực tiểu giao thoa của hai nguồn cùng pha là — phát biểu đúng là: $d_2-d_1=k\lambda$. (Vì): Với hai nguồn cùng pha, cực tiểu khi hiệu đường đi bằng nửa nguyên lần bước sóng. b) Sai: phương án nêu không phù hợp. c) Đúng: Ta có $
\item (Sai) Điều kiện để tại $M$ có cực tiểu giao thoa của hai nguồn cùng pha là — phát biểu $d_2-d_1=(k+0,5)\lambda$ là đúng.
\item (Đúng) Hai nguồn cùng pha, $\lambda=2$ cm. Điểm $M$ có $d_2-d_1=6$ cm là — phát biểu đúng là: cực đại bậc 3.
\item (Sai) Hai nguồn cùng pha, $\lambda=2$ cm. Điểm $M$ có $d_2-d_1=6$ cm là — phát biểu cực đại bậc 2 là đúng. (Vì): _1)/\lambda=6/2=3$, nên $M$ thuộc cực đại bậc 3. d) Sai: phương án nêu không phù hợp
\end{itemchoice}
}
\end{ex}

% ===== Câu 86 =====
% ID: L11C2B12-22-TL
% Mức: VD
\dangbt{Vận dụng giao thoa sóng trong thực tế}
\begin{ex}
% ID: L11C2B12-22-TL
% Mức: VD
Trong dạng “Vận dụng giao thoa sóng trong thực tế”, Hai nguồn cùng pha có $\lambda=2$ cm. Xét điểm $M$ với $d_1=10$ cm, $d_2=15$ cm. Xác định trạng thái giao thoa tại M.
\choiceTF

\loigiai{
Hiệu đường đi $\Delta d=5\,\mathrm{cm}=2,5\lambda$, nên hai sóng ngược pha và $M$ thuộc cực tiểu.
}
\end{ex}

% ===== Câu 87 =====
% ID: L11C2B12-22-TLN
% Mức: VD
\begin{ex}
% ID: L11C2B12-22-TLN
% Mức: VD
Trong dạng “Vận dụng giao thoa sóng trong thực tế”, Sóng có tốc độ $40\,\text{cm}$ /s và tần số $10\,\text{Hz}$. Tính bước sóng theo cm.
\shortans{4}
\loigiai{
$\lambda=v/f=40/10=4$ cm.
}
\end{ex}

% ===== Câu 88 =====
% ID: L11C2B12-22-TN
% Mức: VD
\dangbt{Nhận biết hiện tượng và điều kiện giao thoa sóng}
\begin{ex}
% ID: L11C2B12-22-TN
% Mức: VD
Trong dạng “Nhận biết hiện tượng và điều kiện giao thoa sóng”, Trên đoạn nối hai nguồn cùng pha, khoảng cách giữa hai cực đại liên tiếp bằng
\choice
{\True $\lambda/2$}
{$\lambda$}
{$2\lambda$}
{$\lambda/4$}
\loigiai{
Đáp án A. Trên đoạn nối hai nguồn, hiệu đường đi đổi hai lần theo li độ nên hai cực đại liên tiếp cách nhau $\lambda/2$.
}
\end{ex}

% ===== Câu 89 =====
% ID: L11C2B12-23-DS
% Mức: VD
\dangbt{Xác định cực đại, cực tiểu giao thoa}
\begin{ex}
% ID: L11C2B12-23-DS
% Mức: VD
Xét các phát biểu sau trong dạng “Xác định cực đại, cực tiểu giao thoa”.
\choiceTF
{\True Hai nguồn cùng pha, $\lambda=4$ cm. Điểm $N$ có $d_2-d_1=6$ cm là — phát biểu đúng là: cực đại.}
{Hai nguồn cùng pha, $\lambda=4$ cm. Điểm $N$ có $d_2-d_1=6$ cm là — phát biểu cực tiểu là đúng.}
{\True Trên đoạn nối hai nguồn cùng pha, khoảng cách giữa hai cực đại liên tiếp bằng — phát biểu đúng là: $\lambda/2$.}
{Trên đoạn nối hai nguồn cùng pha, khoảng cách giữa hai cực đại liên tiếp bằng — phát biểu $2\lambda$ là đúng.}
\loigiai{
\begin{itemchoice}
\item (Đúng) Hai nguồn cùng pha, $\lambda=4$ cm. Điểm $N$ có $d_2-d_1=6$ cm là — phát biểu đúng là: cực đại. (Vì): $6/4=1,5$, là nửa nguyên nên $N$ thuộc cực tiểu. b) Sai: phương án nêu không phù hợp. c) Đúng: Trên đoạn nối hai nguồn, hiệu đường đi đổi hai lần theo li độ nên hai cực đại liên tiếp cách nhau $\lambda/2$. d) Sai: phương án nêu không phù hợp
\item (Sai) Hai nguồn cùng pha, $\lambda=4$ cm. Điểm $N$ có $d_2-d_1=6$ cm là — phát biểu cực tiểu là đúng.
\item (Đúng) Trên đoạn nối hai nguồn cùng pha, khoảng cách giữa hai cực đại liên tiếp bằng — phát biểu đúng là: $\lambda/2$.
\item (Sai) Trên đoạn nối hai nguồn cùng pha, khoảng cách giữa hai cực đại liên tiếp bằng — phát biểu $2\lambda$ là đúng.
\end{itemchoice}
}
\end{ex}

% ===== Câu 90 =====
% ID: L11C2B12-23-TL
% Mức: VD
\dangbt{Vận dụng giao thoa sóng trong thực tế}
\begin{ex}
% ID: L11C2B12-23-TL
% Mức: VD
Trong dạng “Vận dụng giao thoa sóng trong thực tế”, Trình bày cách xác định bước sóng từ ảnh các đường cực đại trên đoạn nối hai nguồn.
\choiceTF

\loigiai{
Đo khoảng cách $i$ giữa hai cực đại liên tiếp trên đoạn nối nguồn; do $i=\lambda/2$ nên tính $\lambda=2i$.
}
\end{ex}

% ===== Câu 91 =====
% ID: L11C2B12-23-TLN
% Mức: VD
\begin{ex}
% ID: L11C2B12-23-TLN
% Mức: VD
Trong dạng “Vận dụng giao thoa sóng trong thực tế”, Hai nguồn cùng pha cách nhau $12\,\text{cm}$, $\lambda=3$ cm. Giá trị bậc cực đại lớn nhất trên đoạn nối nguồn là bao nhiêu?
\shortans{4}
\loigiai{
$|k|\le a/\lambda=12/3=4$, nên bậc lớn nhất là 4.
}
\end{ex}

% ===== Câu 92 =====
% ID: L11C2B12-23-TN
% Mức: VD
\dangbt{Nhận biết hiện tượng và điều kiện giao thoa sóng}
\begin{ex}
% ID: L11C2B12-23-TN
% Mức: VD
Trong dạng “Nhận biết hiện tượng và điều kiện giao thoa sóng”, Hiện tượng giao thoa là bằng chứng rõ ràng của
\choice
{\True tính chất sóng}
{tính chất hạt}
{sự bảo toàn khối lượng}
{chuyển động thẳng đều}
\loigiai{
Đáp án A. Giao thoa là hiện tượng đặc trưng của sóng.
}
\end{ex}

% ===== Câu 93 =====
% ID: L11C2B12-24-DS
% Mức: VD
\dangbt{Xác định cực đại, cực tiểu giao thoa}
\begin{ex}
% ID: L11C2B12-24-DS
% Mức: VD
Xét các phát biểu sau trong dạng “Xác định cực đại, cực tiểu giao thoa”.
\choiceTF
{\True Hiện tượng giao thoa là bằng chứng rõ ràng của — phát biểu đúng là: tính chất sóng.}
{Hiện tượng giao thoa là bằng chứng rõ ràng của — phát biểu tính chất hạt là đúng.}
{\True Hai nguồn cùng pha cách nhau $10\,\text{cm}$, $\lambda=2$ cm. Số cực đại trên đoạn nối hai nguồn, kể cả hai nguồn nếu thỏa điều kiện, là — phát biểu đúng là: 9.}
{Hai nguồn cùng pha cách nhau $10\,\text{cm}$, $\lambda=2$ cm. Số cực đại trên đoạn nối hai nguồn, kể cả hai nguồn nếu thỏa điều kiện, là — phát biểu 11 là đúng.}
\loigiai{
\begin{itemchoice}
\item (Đúng) Hiện tượng giao thoa là bằng chứng rõ ràng của — phát biểu đúng là: tính chất sóng. (Vì): Giao thoa là hiện tượng đặc trưng của sóng. b) Sai: phương án nêu không phù hợp. c) Đúng: Các bậc nguyên thỏa $|k|\le 5$, nên có 11 giá trị từ -5 đến 5. d) Sai: phương án nêu không phù hợp
\item (Sai) Hiện tượng giao thoa là bằng chứng rõ ràng của — phát biểu tính chất hạt là đúng.
\item (Đúng) Hai nguồn cùng pha cách nhau $10\,\text{cm}$, $\lambda=2$ cm. Số cực đại trên đoạn nối hai nguồn, kể cả hai nguồn nếu thỏa điều kiện, là — phát biểu đúng là: 9.
\item (Sai) Hai nguồn cùng pha cách nhau $10\,\text{cm}$, $\lambda=2$ cm. Số cực đại trên đoạn nối hai nguồn, kể cả hai nguồn nếu thỏa điều kiện, là — phát biểu 11 là đúng.
\end{itemchoice}
}
\end{ex}

% ===== Câu 94 =====
% ID: L11C2B12-24-TL
% Mức: VD
\dangbt{Vận dụng giao thoa sóng trong thực tế}
\begin{ex}
% ID: L11C2B12-24-TL
% Mức: VD
Trong dạng “Vận dụng giao thoa sóng trong thực tế”, Hai nguồn cùng pha cách nhau $16\,\text{cm}$, bước sóng $4\,\text{cm}$. Xác định số cực đại trên đoạn nối hai nguồn.
\choiceTF

\loigiai{
Điều kiện $|k|\le16/4=4$. Có các giá trị $k=-4,-3,\ldots,4$, tổng cộng 9 cực đại.
}
\end{ex}

% ===== Câu 95 =====
% ID: L11C2B12-24-TLN
% Mức: VDC
\begin{ex}
% ID: L11C2B12-24-TLN
% Mức: VDC
Trong dạng “Vận dụng giao thoa sóng trong thực tế”, Tại $M$, $d_1=18$ cm, $d_2=25$ cm. Tính độ lớn hiệu đường đi theo cm.
\shortans{7}
\loigiai{
$\Delta d=|d_2-d_1|=|25-18|=7$ cm.
}
\end{ex}

% ===== Câu 96 =====
% ID: L11C2B12-24-TN
% Mức: VD
\dangbt{Nhận biết hiện tượng và điều kiện giao thoa sóng}
\begin{ex}
% ID: L11C2B12-24-TN
% Mức: VD
Trong dạng “Nhận biết hiện tượng và điều kiện giao thoa sóng”, Hai nguồn cùng pha cách nhau $10\,\text{cm}$, $\lambda=2$ cm. Số cực đại trên đoạn nối hai nguồn, kể cả hai nguồn nếu thỏa điều kiện, là
\choice
{9}
{10}
{\True 11}
{12}
\loigiai{
Đáp án C. Các bậc nguyên thỏa $|k|\le 5$, nên có 11 giá trị từ -5 đến 5.
}
\end{ex}

% ===== Câu 97 =====
% ID: L11C2B12-25-DS
% Mức: VDC
\dangbt{Xác định cực đại, cực tiểu giao thoa}
\begin{ex}
% ID: L11C2B12-25-DS
% Mức: VDC
Xét các phát biểu sau trong dạng “Xác định cực đại, cực tiểu giao thoa”.
\choiceTF
{\True Tại điểm có cực đại giao thoa, hai sóng thành phần đến đó — phát biểu đúng là: cùng pha.}
{Tại điểm có cực đại giao thoa, hai sóng thành phần đến đó — phát biểu ngược pha là đúng.}
{\True Nếu tần số nguồn tăng gấp đôi còn tốc độ truyền sóng không đổi thì bước sóng — phát biểu đúng là: giảm một nửa.}
{Nếu tần số nguồn tăng gấp đôi còn tốc độ truyền sóng không đổi thì bước sóng — phát biểu không đổi là đúng.}
\loigiai{
\begin{itemchoice}
\item (Đúng) Tại điểm có cực đại giao thoa, hai sóng thành phần đến đó — phát biểu đúng là: cùng pha. (Vì): Cực đại xuất hiện khi hai dao động thành phần tăng cường nhau, tức cùng pha. b) Sai: phương án nêu không phù hợp. c) Đúng: Từ $\lambda=v/f$, khi $f$ tăng gấp đôi thì $\lambda$ giảm một nửa. d) Sai: phương án nêu không phù hợp
\item (Sai) Tại điểm có cực đại giao thoa, hai sóng thành phần đến đó — phát biểu ngược pha là đúng.
\item (Đúng) Nếu tần số nguồn tăng gấp đôi còn tốc độ truyền sóng không đổi thì bước sóng — phát biểu đúng là: giảm một nửa.
\item (Sai) Nếu tần số nguồn tăng gấp đôi còn tốc độ truyền sóng không đổi thì bước sóng — phát biểu không đổi là đúng.
\end{itemchoice}
}
\end{ex}

% ===== Câu 98 =====
% ID: L11C2B12-25-TL
% Mức: VDC
\dangbt{Vận dụng giao thoa sóng trong thực tế}
\begin{ex}
% ID: L11C2B12-25-TL
% Mức: VDC
Trong dạng “Vận dụng giao thoa sóng trong thực tế”, Phân biệt cực đại và cực tiểu giao thoa của hai nguồn cùng pha bằng hiệu đường đi.
\choiceTF

\loigiai{
Cực đại: $\Delta d=k\lambda$. Cực tiểu: $\Delta d=(k+1/2)\lambda$, với $k$ nguyên.
}
\end{ex}

% ===== Câu 99 =====
% ID: L11C2B12-25-TLN
% Mức: TH
\dangbt{Xác định cực đại, cực tiểu giao thoa}
\begin{ex}
% ID: L11C2B12-25-TLN
% Mức: TH
Trong dạng “Xác định cực đại, cực tiểu giao thoa”, Hai nguồn cùng pha, $\lambda=3$ cm. Tại $M$, hiệu đường đi là $12\,\text{cm}$. Xác định bậc cực đại.
\shortans{4}
\loigiai{
$k=(d_2-d_1)/\lambda=12/3=4$.
}
\end{ex}

% ===== Câu 100 =====
% ID: L11C2B12-25-TN
% Mức: VDC
\dangbt{Nhận biết hiện tượng và điều kiện giao thoa sóng}
\begin{ex}
% ID: L11C2B12-25-TN
% Mức: VDC
Trong dạng “Nhận biết hiện tượng và điều kiện giao thoa sóng”, Tại điểm có cực đại giao thoa, hai sóng thành phần đến đó
\choice
{\True cùng pha}
{ngược pha}
{vuông pha}
{khác tần số}
\loigiai{
Đáp án A. Cực đại xuất hiện khi hai dao động thành phần tăng cường nhau, tức cùng pha.
}
\end{ex}

% ===== Câu 101 =====
% ID: L11C2B12-26-DS
% Mức: TH
\dangbt{Xác định khoảng cách giữa các vân giao thoa}
\begin{ex}
% ID: L11C2B12-26-DS
% Mức: TH
Xét các phát biểu sau trong dạng “Xác định khoảng cách giữa các vân giao thoa”.
\choiceTF
{\True Hai nguồn kết hợp là hai nguồn dao động — phát biểu đúng là: cùng phương, cùng tần số và có độ lệch pha không đổi.}
{Hai nguồn kết hợp là hai nguồn dao động — phát biểu cùng biên độ nhưng khác tần số là đúng.}
{\True Điều kiện để tại $M$ có cực đại giao thoa của hai nguồn cùng pha là — phát biểu đúng là: $d_2-d_1=k\lambda$.}
{Điều kiện để tại $M$ có cực đại giao thoa của hai nguồn cùng pha là — phát biểu $d_2+d_1=k\lambda$ là đúng.}
\loigiai{
\begin{itemchoice}
\item (Đúng) Hai nguồn kết hợp là hai nguồn dao động — phát biểu đúng là: cùng phương, cùng tần số và có độ lệch pha không đổi. (Vì): Hai nguồn kết hợp phải cùng phương, cùng tần số và có độ lệch pha không đổi. b) Sai: phương án nêu không phù hợp. c) Đúng: Với hai nguồn cùng pha, cực đại khi hiệu đường đi bằng số nguyên lần bước sóng. d) Sai: phương án nêu không phù hợp
\item (Sai) Hai nguồn kết hợp là hai nguồn dao động — phát biểu cùng biên độ nhưng khác tần số là đúng.
\item (Đúng) Điều kiện để tại $M$ có cực đại giao thoa của hai nguồn cùng pha là — phát biểu đúng là: $d_2-d_1=k\lambda$.
\item (Sai) Điều kiện để tại $M$ có cực đại giao thoa của hai nguồn cùng pha là — phát biểu $d_2+d_1=k\lambda$ là đúng.
\end{itemchoice}
}
\end{ex}

% ===== Câu 102 =====
% ID: L11C2B12-26-TL
% Mức: TH
\dangbt{Xác định cực đại, cực tiểu giao thoa}
\begin{ex}
% ID: L11C2B12-26-TL
% Mức: TH
Trong dạng “Xác định cực đại, cực tiểu giao thoa”, Nêu điều kiện để hai nguồn sóng là hai nguồn kết hợp.
\choiceTF

\loigiai{
Hai nguồn dao động cùng phương, cùng tần số và có độ lệch pha không đổi theo thời gian.
}
\end{ex}

% ===== Câu 103 =====
% ID: L11C2B12-26-TLN
% Mức: TH
\begin{ex}
% ID: L11C2B12-26-TLN
% Mức: TH
Trong dạng “Xác định cực đại, cực tiểu giao thoa”, Hai nguồn cùng pha, $\lambda=4$ cm. Hiệu đường đi tại cực tiểu thứ hai là bao nhiêu theo cm?
\shortans{6}
\loigiai{
Cực tiểu thứ hai ứng với $|d_2-d_1|=(1+0,5)\lambda=1,5\cdot4=6$ cm.
}
\end{ex}

% ===== Câu 104 =====
% ID: L11C2B12-26-TN
% Mức: VDC
\dangbt{Nhận biết hiện tượng và điều kiện giao thoa sóng}
\begin{ex}
% ID: L11C2B12-26-TN
% Mức: VDC
Trong dạng “Nhận biết hiện tượng và điều kiện giao thoa sóng”, Nếu tần số nguồn tăng gấp đôi còn tốc độ truyền sóng không đổi thì bước sóng
\choice
{\True giảm một nửa}
{tăng gấp đôi}
{không đổi}
{tăng bốn lần}
\loigiai{
Đáp án A. Từ $\lambda=v/f$, khi $f$ tăng gấp đôi thì $\lambda$ giảm một nửa.
}
\end{ex}

% ===== Câu 105 =====
% ID: L11C2B12-27-DS
% Mức: TH
\dangbt{Xác định khoảng cách giữa các vân giao thoa}
\begin{ex}
% ID: L11C2B12-27-DS
% Mức: TH
Xét các phát biểu sau trong dạng “Xác định khoảng cách giữa các vân giao thoa”.
\choiceTF
{\True Điều kiện để tại $M$ có cực tiểu giao thoa của hai nguồn cùng pha là — phát biểu đúng là: $d_2-d_1=k\lambda$.}
{Điều kiện để tại $M$ có cực tiểu giao thoa của hai nguồn cùng pha là — phát biểu $d_2-d_1=(k+0,5)\lambda$ là đúng.}
{\True Hai nguồn cùng pha, $\lambda=2$ cm. Điểm $M$ có $d_2-d_1=6$ cm là — phát biểu đúng là: cực đại bậc 3.}
{Hai nguồn cùng pha, $\lambda=2$ cm. Điểm $M$ có $d_2-d_1=6$ cm là — phát biểu cực đại bậc 2 là đúng.}
\loigiai{
\begin{itemchoice}
\item (Đúng) Điều kiện để tại $M$ có cực tiểu giao thoa của hai nguồn cùng pha là — phát biểu đúng là: $d_2-d_1=k\lambda$. (Vì): Với hai nguồn cùng pha, cực tiểu khi hiệu đường đi bằng nửa nguyên lần bước sóng. b) Sai: phương án nêu không phù hợp. c) Đúng: Ta có $
\item (Sai) Điều kiện để tại $M$ có cực tiểu giao thoa của hai nguồn cùng pha là — phát biểu $d_2-d_1=(k+0,5)\lambda$ là đúng.
\item (Đúng) Hai nguồn cùng pha, $\lambda=2$ cm. Điểm $M$ có $d_2-d_1=6$ cm là — phát biểu đúng là: cực đại bậc 3.
\item (Sai) Hai nguồn cùng pha, $\lambda=2$ cm. Điểm $M$ có $d_2-d_1=6$ cm là — phát biểu cực đại bậc 2 là đúng. (Vì): _1)/\lambda=6/2=3$, nên $M$ thuộc cực đại bậc 3. d) Sai: phương án nêu không phù hợp
\end{itemchoice}
}
\end{ex}

% ===== Câu 106 =====
% ID: L11C2B12-27-TL
% Mức: TH
\dangbt{Xác định cực đại, cực tiểu giao thoa}
\begin{ex}
% ID: L11C2B12-27-TL
% Mức: TH
Trong dạng “Xác định cực đại, cực tiểu giao thoa”, Giải thích vì sao ở cực đại giao thoa biên độ dao động lớn.
\choiceTF

\loigiai{
Tại cực đại, hai sóng đến cùng pha nên li độ tổng hợp được cộng đại số và biên độ đạt giá trị lớn.
}
\end{ex}

% ===== Câu 107 =====
% ID: L11C2B12-27-TLN
% Mức: VD
\begin{ex}
% ID: L11C2B12-27-TLN
% Mức: VD
Trong dạng “Xác định cực đại, cực tiểu giao thoa”, Trên đoạn nối hai nguồn, hai cực đại liên tiếp cách nhau $2,5\,\text{cm}$. Tính bước sóng theo cm.
\shortans{5}
\loigiai{
Khoảng cách hai cực đại liên tiếp trên đoạn nối nguồn bằng $\lambda/2$, nên $\lambda=5$ cm.
}
\end{ex}

% ===== Câu 108 =====
% ID: L11C2B12-27-TN
% Mức: NB
\dangbt{Tính hiệu đường đi và bậc giao thoa}
\begin{ex}
% ID: L11C2B12-27-TN
% Mức: NB
Trong dạng “Tính hiệu đường đi và bậc giao thoa”, Hai nguồn kết hợp là hai nguồn dao động
\choice
{\True cùng phương, cùng tần số và có độ lệch pha không đổi}
{cùng biên độ nhưng khác tần số}
{cùng pha ở mọi thời điểm dù khác tần số}
{vuông pha và khác phương}
\loigiai{
Đáp án A. Hai nguồn kết hợp phải cùng phương, cùng tần số và có độ lệch pha không đổi.
}
\end{ex}

% ===== Câu 109 =====
% ID: L11C2B12-28-DS
% Mức: VD
\dangbt{Xác định khoảng cách giữa các vân giao thoa}
\begin{ex}
% ID: L11C2B12-28-DS
% Mức: VD
Xét các phát biểu sau trong dạng “Xác định khoảng cách giữa các vân giao thoa”.
\choiceTF
{\True Hai nguồn cùng pha, $\lambda=4$ cm. Điểm $N$ có $d_2-d_1=6$ cm là — phát biểu đúng là: cực đại.}
{Hai nguồn cùng pha, $\lambda=4$ cm. Điểm $N$ có $d_2-d_1=6$ cm là — phát biểu cực tiểu là đúng.}
{\True Trên đoạn nối hai nguồn cùng pha, khoảng cách giữa hai cực đại liên tiếp bằng — phát biểu đúng là: $\lambda/2$.}
{Trên đoạn nối hai nguồn cùng pha, khoảng cách giữa hai cực đại liên tiếp bằng — phát biểu $2\lambda$ là đúng.}
\loigiai{
\begin{itemchoice}
\item (Đúng) Hai nguồn cùng pha, $\lambda=4$ cm. Điểm $N$ có $d_2-d_1=6$ cm là — phát biểu đúng là: cực đại. (Vì): $6/4=1,5$, là nửa nguyên nên $N$ thuộc cực tiểu. b) Sai: phương án nêu không phù hợp. c) Đúng: Trên đoạn nối hai nguồn, hiệu đường đi đổi hai lần theo li độ nên hai cực đại liên tiếp cách nhau $\lambda/2$. d) Sai: phương án nêu không phù hợp
\item (Sai) Hai nguồn cùng pha, $\lambda=4$ cm. Điểm $N$ có $d_2-d_1=6$ cm là — phát biểu cực tiểu là đúng.
\item (Đúng) Trên đoạn nối hai nguồn cùng pha, khoảng cách giữa hai cực đại liên tiếp bằng — phát biểu đúng là: $\lambda/2$.
\item (Sai) Trên đoạn nối hai nguồn cùng pha, khoảng cách giữa hai cực đại liên tiếp bằng — phát biểu $2\lambda$ là đúng.
\end{itemchoice}
}
\end{ex}

% ===== Câu 110 =====
% ID: L11C2B12-28-TL
% Mức: VD
\dangbt{Xác định cực đại, cực tiểu giao thoa}
\begin{ex}
% ID: L11C2B12-28-TL
% Mức: VD
Trong dạng “Xác định cực đại, cực tiểu giao thoa”, Hai nguồn cùng pha có $\lambda=2$ cm. Xét điểm $M$ với $d_1=10$ cm, $d_2=15$ cm. Xác định trạng thái giao thoa tại M.
\choiceTF

\loigiai{
Hiệu đường đi $\Delta d=5\,\mathrm{cm}=2,5\lambda$, nên hai sóng ngược pha và $M$ thuộc cực tiểu.
}
\end{ex}

% ===== Câu 111 =====
% ID: L11C2B12-28-TLN
% Mức: VD
\begin{ex}
% ID: L11C2B12-28-TLN
% Mức: VD
Trong dạng “Xác định cực đại, cực tiểu giao thoa”, Sóng có tốc độ $40\,\text{cm}$ /s và tần số $10\,\text{Hz}$. Tính bước sóng theo cm.
\shortans{4}
\loigiai{
$\lambda=v/f=40/10=4$ cm.
}
\end{ex}

% ===== Câu 112 =====
% ID: L11C2B12-28-TN
% Mức: NB
\dangbt{Tính hiệu đường đi và bậc giao thoa}
\begin{ex}
% ID: L11C2B12-28-TN
% Mức: NB
Trong dạng “Tính hiệu đường đi và bậc giao thoa”, Điều kiện để tại $M$ có cực đại giao thoa của hai nguồn cùng pha là
\choice
{\True $d_2-d_1=k\lambda$}
{$d_2-d_1=(k+0,5)\lambda$}
{$d_2+d_1=k\lambda$}
{$d_2-d_1=\lambda/4$}
\loigiai{
Đáp án A. Với hai nguồn cùng pha, cực đại khi hiệu đường đi bằng số nguyên lần bước sóng.
}
\end{ex}

% ===== Câu 113 =====
% ID: L11C2B12-29-DS
% Mức: VD
\dangbt{Xác định khoảng cách giữa các vân giao thoa}
\begin{ex}
% ID: L11C2B12-29-DS
% Mức: VD
Xét các phát biểu sau trong dạng “Xác định khoảng cách giữa các vân giao thoa”.
\choiceTF
{\True Hiện tượng giao thoa là bằng chứng rõ ràng của — phát biểu đúng là: tính chất sóng.}
{Hiện tượng giao thoa là bằng chứng rõ ràng của — phát biểu tính chất hạt là đúng.}
{\True Hai nguồn cùng pha cách nhau $10\,\text{cm}$, $\lambda=2$ cm. Số cực đại trên đoạn nối hai nguồn, kể cả hai nguồn nếu thỏa điều kiện, là — phát biểu đúng là: 9.}
{Hai nguồn cùng pha cách nhau $10\,\text{cm}$, $\lambda=2$ cm. Số cực đại trên đoạn nối hai nguồn, kể cả hai nguồn nếu thỏa điều kiện, là — phát biểu 11 là đúng.}
\loigiai{
\begin{itemchoice}
\item (Đúng) Hiện tượng giao thoa là bằng chứng rõ ràng của — phát biểu đúng là: tính chất sóng. (Vì): Giao thoa là hiện tượng đặc trưng của sóng. b) Sai: phương án nêu không phù hợp. c) Đúng: Các bậc nguyên thỏa $|k|\le 5$, nên có 11 giá trị từ -5 đến 5. d) Sai: phương án nêu không phù hợp
\item (Sai) Hiện tượng giao thoa là bằng chứng rõ ràng của — phát biểu tính chất hạt là đúng.
\item (Đúng) Hai nguồn cùng pha cách nhau $10\,\text{cm}$, $\lambda=2$ cm. Số cực đại trên đoạn nối hai nguồn, kể cả hai nguồn nếu thỏa điều kiện, là — phát biểu đúng là: 9.
\item (Sai) Hai nguồn cùng pha cách nhau $10\,\text{cm}$, $\lambda=2$ cm. Số cực đại trên đoạn nối hai nguồn, kể cả hai nguồn nếu thỏa điều kiện, là — phát biểu 11 là đúng.
\end{itemchoice}
}
\end{ex}

% ===== Câu 114 =====
% ID: L11C2B12-29-TL
% Mức: VD
\dangbt{Xác định cực đại, cực tiểu giao thoa}
\begin{ex}
% ID: L11C2B12-29-TL
% Mức: VD
Trong dạng “Xác định cực đại, cực tiểu giao thoa”, Trình bày cách xác định bước sóng từ ảnh các đường cực đại trên đoạn nối hai nguồn.
\choiceTF

\loigiai{
Đo khoảng cách $i$ giữa hai cực đại liên tiếp trên đoạn nối nguồn; do $i=\lambda/2$ nên tính $\lambda=2i$.
}
\end{ex}

% ===== Câu 115 =====
% ID: L11C2B12-29-TLN
% Mức: VD
\begin{ex}
% ID: L11C2B12-29-TLN
% Mức: VD
Trong dạng “Xác định cực đại, cực tiểu giao thoa”, Hai nguồn cùng pha cách nhau $12\,\text{cm}$, $\lambda=3$ cm. Giá trị bậc cực đại lớn nhất trên đoạn nối nguồn là bao nhiêu?
\shortans{4}
\loigiai{
$|k|\le a/\lambda=12/3=4$, nên bậc lớn nhất là 4.
}
\end{ex}

% ===== Câu 116 =====
% ID: L11C2B12-29-TN
% Mức: TH
\dangbt{Tính hiệu đường đi và bậc giao thoa}
\begin{ex}
% ID: L11C2B12-29-TN
% Mức: TH
Trong dạng “Tính hiệu đường đi và bậc giao thoa”, Điều kiện để tại $M$ có cực tiểu giao thoa của hai nguồn cùng pha là
\choice
{$d_2-d_1=k\lambda$}
{\True $d_2-d_1=(k+0,5)\lambda$}
{$d_2+d_1=(k+0,5)\lambda$}
{$d_2-d_1=2k\lambda$}
\loigiai{
Đáp án B. Với hai nguồn cùng pha, cực tiểu khi hiệu đường đi bằng nửa nguyên lần bước sóng.
}
\end{ex}

% ===== Câu 117 =====
% ID: L11C2B12-30-DS
% Mức: VDC
\dangbt{Xác định khoảng cách giữa các vân giao thoa}
\begin{ex}
% ID: L11C2B12-30-DS
% Mức: VDC
Xét các phát biểu sau trong dạng “Xác định khoảng cách giữa các vân giao thoa”.
\choiceTF
{\True Tại điểm có cực đại giao thoa, hai sóng thành phần đến đó — phát biểu đúng là: cùng pha.}
{Tại điểm có cực đại giao thoa, hai sóng thành phần đến đó — phát biểu ngược pha là đúng.}
{\True Nếu tần số nguồn tăng gấp đôi còn tốc độ truyền sóng không đổi thì bước sóng — phát biểu đúng là: giảm một nửa.}
{Nếu tần số nguồn tăng gấp đôi còn tốc độ truyền sóng không đổi thì bước sóng — phát biểu không đổi là đúng.}
\loigiai{
\begin{itemchoice}
\item (Đúng) Tại điểm có cực đại giao thoa, hai sóng thành phần đến đó — phát biểu đúng là: cùng pha. (Vì): Cực đại xuất hiện khi hai dao động thành phần tăng cường nhau, tức cùng pha. b) Sai: phương án nêu không phù hợp. c) Đúng: Từ $\lambda=v/f$, khi $f$ tăng gấp đôi thì $\lambda$ giảm một nửa. d) Sai: phương án nêu không phù hợp
\item (Sai) Tại điểm có cực đại giao thoa, hai sóng thành phần đến đó — phát biểu ngược pha là đúng.
\item (Đúng) Nếu tần số nguồn tăng gấp đôi còn tốc độ truyền sóng không đổi thì bước sóng — phát biểu đúng là: giảm một nửa.
\item (Sai) Nếu tần số nguồn tăng gấp đôi còn tốc độ truyền sóng không đổi thì bước sóng — phát biểu không đổi là đúng.
\end{itemchoice}
}
\end{ex}

% ===== Câu 118 =====
% ID: L11C2B12-30-TL
% Mức: VD
\dangbt{Xác định cực đại, cực tiểu giao thoa}
\begin{ex}
% ID: L11C2B12-30-TL
% Mức: VD
Trong dạng “Xác định cực đại, cực tiểu giao thoa”, Hai nguồn cùng pha cách nhau $16\,\text{cm}$, bước sóng $4\,\text{cm}$. Xác định số cực đại trên đoạn nối hai nguồn.
\choiceTF

\loigiai{
Điều kiện $|k|\le16/4=4$. Có các giá trị $k=-4,-3,\ldots,4$, tổng cộng 9 cực đại.
}
\end{ex}

% ===== Câu 119 =====
% ID: L11C2B12-30-TLN
% Mức: VDC
\begin{ex}
% ID: L11C2B12-30-TLN
% Mức: VDC
Trong dạng “Xác định cực đại, cực tiểu giao thoa”, Tại $M$, $d_1=18$ cm, $d_2=25$ cm. Tính độ lớn hiệu đường đi theo cm.
\shortans{7}
\loigiai{
$\Delta d=|d_2-d_1|=|25-18|=7$ cm.
}
\end{ex}

% ===== Câu 120 =====
% ID: L11C2B12-30-TN
% Mức: TH
\dangbt{Tính hiệu đường đi và bậc giao thoa}
\begin{ex}
% ID: L11C2B12-30-TN
% Mức: TH
Trong dạng “Tính hiệu đường đi và bậc giao thoa”, Hai nguồn cùng pha, $\lambda=2$ cm. Điểm $M$ có $d_2-d_1=6$ cm là
\choice
{\True cực đại bậc 3}
{cực tiểu thứ 3}
{cực đại bậc 2}
{không giao thoa}
\loigiai{
Đáp án A. Ta có $(d_2-d_1)/\lambda=6/2=3$, nên $M$ thuộc cực đại bậc 3.
}
\end{ex}

% ===== Câu 121 =====
% ID: L11C2B12-31-TL
% Mức: VDC
\dangbt{Xác định cực đại, cực tiểu giao thoa}
\begin{ex}
% ID: L11C2B12-31-TL
% Mức: VDC
Trong dạng “Xác định cực đại, cực tiểu giao thoa”, Phân biệt cực đại và cực tiểu giao thoa của hai nguồn cùng pha bằng hiệu đường đi.
\choiceTF

\loigiai{
Cực đại: $\Delta d=k\lambda$. Cực tiểu: $\Delta d=(k+1/2)\lambda$, với $k$ nguyên.
}
\end{ex}

% ===== Câu 122 =====
% ID: L11C2B12-31-TLN
% Mức: TH
\dangbt{Xác định khoảng cách giữa các vân giao thoa}
\begin{ex}
% ID: L11C2B12-31-TLN
% Mức: TH
Trong dạng “Xác định khoảng cách giữa các vân giao thoa”, Hai nguồn cùng pha, $\lambda=3$ cm. Tại $M$, hiệu đường đi là $12\,\text{cm}$. Xác định bậc cực đại.
\shortans{4}
\loigiai{
$k=(d_2-d_1)/\lambda=12/3=4$.
}
\end{ex}

% ===== Câu 123 =====
% ID: L11C2B12-31-TN
% Mức: TH
\dangbt{Tính hiệu đường đi và bậc giao thoa}
\begin{ex}
% ID: L11C2B12-31-TN
% Mức: TH
Trong dạng “Tính hiệu đường đi và bậc giao thoa”, Hai nguồn cùng pha, $\lambda=4$ cm. Điểm $N$ có $d_2-d_1=6$ cm là
\choice
{cực đại}
{\True cực tiểu}
{điểm đứng yên do không có sóng tới}
{không xác định}
\loigiai{
Đáp án B. $6/4=1,5$, là nửa nguyên nên $N$ thuộc cực tiểu.
}
\end{ex}

% ===== Câu 124 =====
% ID: L11C2B12-32-TL
% Mức: TH
\dangbt{Xác định khoảng cách giữa các vân giao thoa}
\begin{ex}
% ID: L11C2B12-32-TL
% Mức: TH
Trong dạng “Xác định khoảng cách giữa các vân giao thoa”, Nêu điều kiện để hai nguồn sóng là hai nguồn kết hợp.
\choiceTF

\loigiai{
Hai nguồn dao động cùng phương, cùng tần số và có độ lệch pha không đổi theo thời gian.
}
\end{ex}

% ===== Câu 125 =====
% ID: L11C2B12-32-TLN
% Mức: TH
\begin{ex}
% ID: L11C2B12-32-TLN
% Mức: TH
Trong dạng “Xác định khoảng cách giữa các vân giao thoa”, Hai nguồn cùng pha, $\lambda=4$ cm. Hiệu đường đi tại cực tiểu thứ hai là bao nhiêu theo cm?
\shortans{6}
\loigiai{
Cực tiểu thứ hai ứng với $|d_2-d_1|=(1+0,5)\lambda=1,5\cdot4=6$ cm.
}
\end{ex}

% ===== Câu 126 =====
% ID: L11C2B12-32-TN
% Mức: VD
\dangbt{Tính hiệu đường đi và bậc giao thoa}
\begin{ex}
% ID: L11C2B12-32-TN
% Mức: VD
Trong dạng “Tính hiệu đường đi và bậc giao thoa”, Trên đoạn nối hai nguồn cùng pha, khoảng cách giữa hai cực đại liên tiếp bằng
\choice
{\True $\lambda/2$}
{$\lambda$}
{$2\lambda$}
{$\lambda/4$}
\loigiai{
Đáp án A. Trên đoạn nối hai nguồn, hiệu đường đi đổi hai lần theo li độ nên hai cực đại liên tiếp cách nhau $\lambda/2$.
}
\end{ex}

% ===== Câu 127 =====
% ID: L11C2B12-33-TL
% Mức: TH
\dangbt{Xác định khoảng cách giữa các vân giao thoa}
\begin{ex}
% ID: L11C2B12-33-TL
% Mức: TH
Trong dạng “Xác định khoảng cách giữa các vân giao thoa”, Giải thích vì sao ở cực đại giao thoa biên độ dao động lớn.
\choiceTF

\loigiai{
Tại cực đại, hai sóng đến cùng pha nên li độ tổng hợp được cộng đại số và biên độ đạt giá trị lớn.
}
\end{ex}

% ===== Câu 128 =====
% ID: L11C2B12-33-TLN
% Mức: VD
\begin{ex}
% ID: L11C2B12-33-TLN
% Mức: VD
Trong dạng “Xác định khoảng cách giữa các vân giao thoa”, Trên đoạn nối hai nguồn, hai cực đại liên tiếp cách nhau $2,5\,\text{cm}$. Tính bước sóng theo cm.
\shortans{5}
\loigiai{
Khoảng cách hai cực đại liên tiếp trên đoạn nối nguồn bằng $\lambda/2$, nên $\lambda=5$ cm.
}
\end{ex}

% ===== Câu 129 =====
% ID: L11C2B12-33-TN
% Mức: VD
\dangbt{Tính hiệu đường đi và bậc giao thoa}
\begin{ex}
% ID: L11C2B12-33-TN
% Mức: VD
Trong dạng “Tính hiệu đường đi và bậc giao thoa”, Hiện tượng giao thoa là bằng chứng rõ ràng của
\choice
{\True tính chất sóng}
{tính chất hạt}
{sự bảo toàn khối lượng}
{chuyển động thẳng đều}
\loigiai{
Đáp án A. Giao thoa là hiện tượng đặc trưng của sóng.
}
\end{ex}

% ===== Câu 130 =====
% ID: L11C2B12-34-TL
% Mức: VD
\dangbt{Xác định khoảng cách giữa các vân giao thoa}
\begin{ex}
% ID: L11C2B12-34-TL
% Mức: VD
Trong dạng “Xác định khoảng cách giữa các vân giao thoa”, Hai nguồn cùng pha có $\lambda=2$ cm. Xét điểm $M$ với $d_1=10$ cm, $d_2=15$ cm. Xác định trạng thái giao thoa tại M.
\choiceTF

\loigiai{
Hiệu đường đi $\Delta d=5\,\mathrm{cm}=2,5\lambda$, nên hai sóng ngược pha và $M$ thuộc cực tiểu.
}
\end{ex}

% ===== Câu 131 =====
% ID: L11C2B12-34-TLN
% Mức: VD
\begin{ex}
% ID: L11C2B12-34-TLN
% Mức: VD
Trong dạng “Xác định khoảng cách giữa các vân giao thoa”, Sóng có tốc độ $40\,\text{cm}$ /s và tần số $10\,\text{Hz}$. Tính bước sóng theo cm.
\shortans{4}
\loigiai{
$\lambda=v/f=40/10=4$ cm.
}
\end{ex}

% ===== Câu 132 =====
% ID: L11C2B12-34-TN
% Mức: VD
\dangbt{Tính hiệu đường đi và bậc giao thoa}
\begin{ex}
% ID: L11C2B12-34-TN
% Mức: VD
Trong dạng “Tính hiệu đường đi và bậc giao thoa”, Hai nguồn cùng pha cách nhau $10\,\text{cm}$, $\lambda=2$ cm. Số cực đại trên đoạn nối hai nguồn, kể cả hai nguồn nếu thỏa điều kiện, là
\choice
{9}
{10}
{\True 11}
{12}
\loigiai{
Đáp án C. Các bậc nguyên thỏa $|k|\le 5$, nên có 11 giá trị từ -5 đến 5.
}
\end{ex}

% ===== Câu 133 =====
% ID: L11C2B12-35-TL
% Mức: VD
\dangbt{Xác định khoảng cách giữa các vân giao thoa}
\begin{ex}
% ID: L11C2B12-35-TL
% Mức: VD
Trong dạng “Xác định khoảng cách giữa các vân giao thoa”, Trình bày cách xác định bước sóng từ ảnh các đường cực đại trên đoạn nối hai nguồn.
\choiceTF

\loigiai{
Đo khoảng cách $i$ giữa hai cực đại liên tiếp trên đoạn nối nguồn; do $i=\lambda/2$ nên tính $\lambda=2i$.
}
\end{ex}

% ===== Câu 134 =====
% ID: L11C2B12-35-TLN
% Mức: VD
\begin{ex}
% ID: L11C2B12-35-TLN
% Mức: VD
Trong dạng “Xác định khoảng cách giữa các vân giao thoa”, Hai nguồn cùng pha cách nhau $12\,\text{cm}$, $\lambda=3$ cm. Giá trị bậc cực đại lớn nhất trên đoạn nối nguồn là bao nhiêu?
\shortans{4}
\loigiai{
$|k|\le a/\lambda=12/3=4$, nên bậc lớn nhất là 4.
}
\end{ex}

% ===== Câu 135 =====
% ID: L11C2B12-35-TN
% Mức: VDC
\dangbt{Tính hiệu đường đi và bậc giao thoa}
\begin{ex}
% ID: L11C2B12-35-TN
% Mức: VDC
Trong dạng “Tính hiệu đường đi và bậc giao thoa”, Tại điểm có cực đại giao thoa, hai sóng thành phần đến đó
\choice
{\True cùng pha}
{ngược pha}
{vuông pha}
{khác tần số}
\loigiai{
Đáp án A. Cực đại xuất hiện khi hai dao động thành phần tăng cường nhau, tức cùng pha.
}
\end{ex}

% ===== Câu 136 =====
% ID: L11C2B12-36-TL
% Mức: VD
\dangbt{Xác định khoảng cách giữa các vân giao thoa}
\begin{ex}
% ID: L11C2B12-36-TL
% Mức: VD
Trong dạng “Xác định khoảng cách giữa các vân giao thoa”, Hai nguồn cùng pha cách nhau $16\,\text{cm}$, bước sóng $4\,\text{cm}$. Xác định số cực đại trên đoạn nối hai nguồn.
\choiceTF

\loigiai{
Điều kiện $|k|\le16/4=4$. Có các giá trị $k=-4,-3,\ldots,4$, tổng cộng 9 cực đại.
}
\end{ex}

% ===== Câu 137 =====
% ID: L11C2B12-36-TLN
% Mức: VDC
\begin{ex}
% ID: L11C2B12-36-TLN
% Mức: VDC
Trong dạng “Xác định khoảng cách giữa các vân giao thoa”, Tại $M$, $d_1=18$ cm, $d_2=25$ cm. Tính độ lớn hiệu đường đi theo cm.
\shortans{7}
\loigiai{
$\Delta d=|d_2-d_1|=|25-18|=7$ cm.
}
\end{ex}

% ===== Câu 138 =====
% ID: L11C2B12-36-TN
% Mức: VDC
\dangbt{Tính hiệu đường đi và bậc giao thoa}
\begin{ex}
% ID: L11C2B12-36-TN
% Mức: VDC
Trong dạng “Tính hiệu đường đi và bậc giao thoa”, Nếu tần số nguồn tăng gấp đôi còn tốc độ truyền sóng không đổi thì bước sóng
\choice
{\True giảm một nửa}
{tăng gấp đôi}
{không đổi}
{tăng bốn lần}
\loigiai{
Đáp án A. Từ $\lambda=v/f$, khi $f$ tăng gấp đôi thì $\lambda$ giảm một nửa.
}
\end{ex}

% ===== Câu 139 =====
% ID: L11C2B12-37-TL
% Mức: VDC
\dangbt{Xác định khoảng cách giữa các vân giao thoa}
\begin{ex}
% ID: L11C2B12-37-TL
% Mức: VDC
Trong dạng “Xác định khoảng cách giữa các vân giao thoa”, Phân biệt cực đại và cực tiểu giao thoa của hai nguồn cùng pha bằng hiệu đường đi.
\choiceTF

\loigiai{
Cực đại: $\Delta d=k\lambda$. Cực tiểu: $\Delta d=(k+1/2)\lambda$, với $k$ nguyên.
}
\end{ex}

% ===== Câu 140 =====
% ID: L11C2B12-37-TN
% Mức: NB
\dangbt{Vận dụng giao thoa sóng trong thực tế}
\begin{ex}
% ID: L11C2B12-37-TN
% Mức: NB
Trong dạng “Vận dụng giao thoa sóng trong thực tế”, Hai nguồn kết hợp là hai nguồn dao động
\choice
{\True cùng phương, cùng tần số và có độ lệch pha không đổi}
{cùng biên độ nhưng khác tần số}
{cùng pha ở mọi thời điểm dù khác tần số}
{vuông pha và khác phương}
\loigiai{
Đáp án A. Hai nguồn kết hợp phải cùng phương, cùng tần số và có độ lệch pha không đổi.
}
\end{ex}

% ===== Câu 141 =====
% ID: L11C2B12-38-TN
% Mức: NB
\begin{ex}
% ID: L11C2B12-38-TN
% Mức: NB
Trong dạng “Vận dụng giao thoa sóng trong thực tế”, Điều kiện để tại $M$ có cực đại giao thoa của hai nguồn cùng pha là
\choice
{\True $d_2-d_1=k\lambda$}
{$d_2-d_1=(k+0,5)\lambda$}
{$d_2+d_1=k\lambda$}
{$d_2-d_1=\lambda/4$}
\loigiai{
Đáp án A. Với hai nguồn cùng pha, cực đại khi hiệu đường đi bằng số nguyên lần bước sóng.
}
\end{ex}

% ===== Câu 142 =====
% ID: L11C2B12-39-TN
% Mức: TH
\begin{ex}
% ID: L11C2B12-39-TN
% Mức: TH
Trong dạng “Vận dụng giao thoa sóng trong thực tế”, Điều kiện để tại $M$ có cực tiểu giao thoa của hai nguồn cùng pha là
\choice
{$d_2-d_1=k\lambda$}
{\True $d_2-d_1=(k+0,5)\lambda$}
{$d_2+d_1=(k+0,5)\lambda$}
{$d_2-d_1=2k\lambda$}
\loigiai{
Đáp án B. Với hai nguồn cùng pha, cực tiểu khi hiệu đường đi bằng nửa nguyên lần bước sóng.
}
\end{ex}

% ===== Câu 143 =====
% ID: L11C2B12-40-TN
% Mức: TH
\begin{ex}
% ID: L11C2B12-40-TN
% Mức: TH
Trong dạng “Vận dụng giao thoa sóng trong thực tế”, Hai nguồn cùng pha, $\lambda=2$ cm. Điểm $M$ có $d_2-d_1=6$ cm là
\choice
{\True cực đại bậc 3}
{cực tiểu thứ 3}
{cực đại bậc 2}
{không giao thoa}
\loigiai{
Đáp án A. Ta có $(d_2-d_1)/\lambda=6/2=3$, nên $M$ thuộc cực đại bậc 3.
}
\end{ex}

% ===== Câu 144 =====
% ID: L11C2B12-41-TN
% Mức: TH
\begin{ex}
% ID: L11C2B12-41-TN
% Mức: TH
Trong dạng “Vận dụng giao thoa sóng trong thực tế”, Hai nguồn cùng pha, $\lambda=4$ cm. Điểm $N$ có $d_2-d_1=6$ cm là
\choice
{cực đại}
{\True cực tiểu}
{điểm đứng yên do không có sóng tới}
{không xác định}
\loigiai{
Đáp án B. $6/4=1,5$, là nửa nguyên nên $N$ thuộc cực tiểu.
}
\end{ex}

% ===== Câu 145 =====
% ID: L11C2B12-42-TN
% Mức: VD
\begin{ex}
% ID: L11C2B12-42-TN
% Mức: VD
Trong dạng “Vận dụng giao thoa sóng trong thực tế”, Trên đoạn nối hai nguồn cùng pha, khoảng cách giữa hai cực đại liên tiếp bằng
\choice
{\True $\lambda/2$}
{$\lambda$}
{$2\lambda$}
{$\lambda/4$}
\loigiai{
Đáp án A. Trên đoạn nối hai nguồn, hiệu đường đi đổi hai lần theo li độ nên hai cực đại liên tiếp cách nhau $\lambda/2$.
}
\end{ex}

% ===== Câu 146 =====
% ID: L11C2B12-43-TN
% Mức: VD
\begin{ex}
% ID: L11C2B12-43-TN
% Mức: VD
Trong dạng “Vận dụng giao thoa sóng trong thực tế”, Hiện tượng giao thoa là bằng chứng rõ ràng của
\choice
{\True tính chất sóng}
{tính chất hạt}
{sự bảo toàn khối lượng}
{chuyển động thẳng đều}
\loigiai{
Đáp án A. Giao thoa là hiện tượng đặc trưng của sóng.
}
\end{ex}

% ===== Câu 147 =====
% ID: L11C2B12-44-TN
% Mức: VD
\begin{ex}
% ID: L11C2B12-44-TN
% Mức: VD
Trong dạng “Vận dụng giao thoa sóng trong thực tế”, Hai nguồn cùng pha cách nhau $10\,\text{cm}$, $\lambda=2$ cm. Số cực đại trên đoạn nối hai nguồn, kể cả hai nguồn nếu thỏa điều kiện, là
\choice
{9}
{10}
{\True 11}
{12}
\loigiai{
Đáp án C. Các bậc nguyên thỏa $|k|\le 5$, nên có 11 giá trị từ -5 đến 5.
}
\end{ex}

% ===== Câu 148 =====
% ID: L11C2B12-45-TN
% Mức: VDC
\begin{ex}
% ID: L11C2B12-45-TN
% Mức: VDC
Trong dạng “Vận dụng giao thoa sóng trong thực tế”, Tại điểm có cực đại giao thoa, hai sóng thành phần đến đó
\choice
{\True cùng pha}
{ngược pha}
{vuông pha}
{khác tần số}
\loigiai{
Đáp án A. Cực đại xuất hiện khi hai dao động thành phần tăng cường nhau, tức cùng pha.
}
\end{ex}

% ===== Câu 149 =====
% ID: L11C2B12-46-TN
% Mức: VDC
\begin{ex}
% ID: L11C2B12-46-TN
% Mức: VDC
Trong dạng “Vận dụng giao thoa sóng trong thực tế”, Nếu tần số nguồn tăng gấp đôi còn tốc độ truyền sóng không đổi thì bước sóng
\choice
{\True giảm một nửa}
{tăng gấp đôi}
{không đổi}
{tăng bốn lần}
\loigiai{
Đáp án A. Từ $\lambda=v/f$, khi $f$ tăng gấp đôi thì $\lambda$ giảm một nửa.
}
\end{ex}

% ===== Câu 150 =====
% ID: L11C2B12-47-TN
% Mức: NB
\dangbt{Xác định cực đại, cực tiểu giao thoa}
\begin{ex}
% ID: L11C2B12-47-TN
% Mức: NB
Trong dạng “Xác định cực đại, cực tiểu giao thoa”, Hai nguồn kết hợp là hai nguồn dao động
\choice
{\True cùng phương, cùng tần số và có độ lệch pha không đổi}
{cùng biên độ nhưng khác tần số}
{cùng pha ở mọi thời điểm dù khác tần số}
{vuông pha và khác phương}
\loigiai{
Đáp án A. Hai nguồn kết hợp phải cùng phương, cùng tần số và có độ lệch pha không đổi.
}
\end{ex}

% ===== Câu 151 =====
% ID: L11C2B12-48-TN
% Mức: NB
\begin{ex}
% ID: L11C2B12-48-TN
% Mức: NB
Trong dạng “Xác định cực đại, cực tiểu giao thoa”, Điều kiện để tại $M$ có cực đại giao thoa của hai nguồn cùng pha là
\choice
{\True $d_2-d_1=k\lambda$}
{$d_2-d_1=(k+0,5)\lambda$}
{$d_2+d_1=k\lambda$}
{$d_2-d_1=\lambda/4$}
\loigiai{
Đáp án A. Với hai nguồn cùng pha, cực đại khi hiệu đường đi bằng số nguyên lần bước sóng.
}
\end{ex}

% ===== Câu 152 =====
% ID: L11C2B12-49-TN
% Mức: TH
\begin{ex}
% ID: L11C2B12-49-TN
% Mức: TH
Trong dạng “Xác định cực đại, cực tiểu giao thoa”, Điều kiện để tại $M$ có cực tiểu giao thoa của hai nguồn cùng pha là
\choice
{$d_2-d_1=k\lambda$}
{\True $d_2-d_1=(k+0,5)\lambda$}
{$d_2+d_1=(k+0,5)\lambda$}
{$d_2-d_1=2k\lambda$}
\loigiai{
Đáp án B. Với hai nguồn cùng pha, cực tiểu khi hiệu đường đi bằng nửa nguyên lần bước sóng.
}
\end{ex}

% ===== Câu 153 =====
% ID: L11C2B12-50-TN
% Mức: TH
\begin{ex}
% ID: L11C2B12-50-TN
% Mức: TH
Trong dạng “Xác định cực đại, cực tiểu giao thoa”, Hai nguồn cùng pha, $\lambda=2$ cm. Điểm $M$ có $d_2-d_1=6$ cm là
\choice
{\True cực đại bậc 3}
{cực tiểu thứ 3}
{cực đại bậc 2}
{không giao thoa}
\loigiai{
Đáp án A. Ta có $(d_2-d_1)/\lambda=6/2=3$, nên $M$ thuộc cực đại bậc 3.
}
\end{ex}

% ===== Câu 154 =====
% ID: L11C2B12-51-TN
% Mức: TH
\begin{ex}
% ID: L11C2B12-51-TN
% Mức: TH
Trong dạng “Xác định cực đại, cực tiểu giao thoa”, Hai nguồn cùng pha, $\lambda=4$ cm. Điểm $N$ có $d_2-d_1=6$ cm là
\choice
{cực đại}
{\True cực tiểu}
{điểm đứng yên do không có sóng tới}
{không xác định}
\loigiai{
Đáp án B. $6/4=1,5$, là nửa nguyên nên $N$ thuộc cực tiểu.
}
\end{ex}

% ===== Câu 155 =====
% ID: L11C2B12-52-TN
% Mức: VD
\begin{ex}
% ID: L11C2B12-52-TN
% Mức: VD
Trong dạng “Xác định cực đại, cực tiểu giao thoa”, Trên đoạn nối hai nguồn cùng pha, khoảng cách giữa hai cực đại liên tiếp bằng
\choice
{\True $\lambda/2$}
{$\lambda$}
{$2\lambda$}
{$\lambda/4$}
\loigiai{
Đáp án A. Trên đoạn nối hai nguồn, hiệu đường đi đổi hai lần theo li độ nên hai cực đại liên tiếp cách nhau $\lambda/2$.
}
\end{ex}

% ===== Câu 156 =====
% ID: L11C2B12-53-TN
% Mức: VD
\begin{ex}
% ID: L11C2B12-53-TN
% Mức: VD
Trong dạng “Xác định cực đại, cực tiểu giao thoa”, Hiện tượng giao thoa là bằng chứng rõ ràng của
\choice
{\True tính chất sóng}
{tính chất hạt}
{sự bảo toàn khối lượng}
{chuyển động thẳng đều}
\loigiai{
Đáp án A. Giao thoa là hiện tượng đặc trưng của sóng.
}
\end{ex}

% ===== Câu 157 =====
% ID: L11C2B12-54-TN
% Mức: VD
\begin{ex}
% ID: L11C2B12-54-TN
% Mức: VD
Trong dạng “Xác định cực đại, cực tiểu giao thoa”, Hai nguồn cùng pha cách nhau $10\,\text{cm}$, $\lambda=2$ cm. Số cực đại trên đoạn nối hai nguồn, kể cả hai nguồn nếu thỏa điều kiện, là
\choice
{9}
{10}
{\True 11}
{12}
\loigiai{
Đáp án C. Các bậc nguyên thỏa $|k|\le 5$, nên có 11 giá trị từ -5 đến 5.
}
\end{ex}

% ===== Câu 158 =====
% ID: L11C2B12-55-TN
% Mức: VDC
\begin{ex}
% ID: L11C2B12-55-TN
% Mức: VDC
Trong dạng “Xác định cực đại, cực tiểu giao thoa”, Tại điểm có cực đại giao thoa, hai sóng thành phần đến đó
\choice
{\True cùng pha}
{ngược pha}
{vuông pha}
{khác tần số}
\loigiai{
Đáp án A. Cực đại xuất hiện khi hai dao động thành phần tăng cường nhau, tức cùng pha.
}
\end{ex}

% ===== Câu 159 =====
% ID: L11C2B12-56-TN
% Mức: VDC
\begin{ex}
% ID: L11C2B12-56-TN
% Mức: VDC
Trong dạng “Xác định cực đại, cực tiểu giao thoa”, Nếu tần số nguồn tăng gấp đôi còn tốc độ truyền sóng không đổi thì bước sóng
\choice
{\True giảm một nửa}
{tăng gấp đôi}
{không đổi}
{tăng bốn lần}
\loigiai{
Đáp án A. Từ $\lambda=v/f$, khi $f$ tăng gấp đôi thì $\lambda$ giảm một nửa.
}
\end{ex}

% ===== Câu 160 =====
% ID: L11C2B12-57-TN
% Mức: NB
\dangbt{Xác định khoảng cách giữa các vân giao thoa}
\begin{ex}
% ID: L11C2B12-57-TN
% Mức: NB
Trong dạng “Xác định khoảng cách giữa các vân giao thoa”, Hai nguồn kết hợp là hai nguồn dao động
\choice
{\True cùng phương, cùng tần số và có độ lệch pha không đổi}
{cùng biên độ nhưng khác tần số}
{cùng pha ở mọi thời điểm dù khác tần số}
{vuông pha và khác phương}
\loigiai{
Đáp án A. Hai nguồn kết hợp phải cùng phương, cùng tần số và có độ lệch pha không đổi.
}
\end{ex}

% ===== Câu 161 =====
% ID: L11C2B12-58-TN
% Mức: NB
\begin{ex}
% ID: L11C2B12-58-TN
% Mức: NB
Trong dạng “Xác định khoảng cách giữa các vân giao thoa”, Điều kiện để tại $M$ có cực đại giao thoa của hai nguồn cùng pha là
\choice
{\True $d_2-d_1=k\lambda$}
{$d_2-d_1=(k+0,5)\lambda$}
{$d_2+d_1=k\lambda$}
{$d_2-d_1=\lambda/4$}
\loigiai{
Đáp án A. Với hai nguồn cùng pha, cực đại khi hiệu đường đi bằng số nguyên lần bước sóng.
}
\end{ex}

% ===== Câu 162 =====
% ID: L11C2B12-59-TN
% Mức: TH
\begin{ex}
% ID: L11C2B12-59-TN
% Mức: TH
Trong dạng “Xác định khoảng cách giữa các vân giao thoa”, Điều kiện để tại $M$ có cực tiểu giao thoa của hai nguồn cùng pha là
\choice
{$d_2-d_1=k\lambda$}
{\True $d_2-d_1=(k+0,5)\lambda$}
{$d_2+d_1=(k+0,5)\lambda$}
{$d_2-d_1=2k\lambda$}
\loigiai{
Đáp án B. Với hai nguồn cùng pha, cực tiểu khi hiệu đường đi bằng nửa nguyên lần bước sóng.
}
\end{ex}

% ===== Câu 163 =====
% ID: L11C2B12-60-TN
% Mức: TH
\begin{ex}
% ID: L11C2B12-60-TN
% Mức: TH
Trong dạng “Xác định khoảng cách giữa các vân giao thoa”, Hai nguồn cùng pha, $\lambda=2$ cm. Điểm $M$ có $d_2-d_1=6$ cm là
\choice
{\True cực đại bậc 3}
{cực tiểu thứ 3}
{cực đại bậc 2}
{không giao thoa}
\loigiai{
Đáp án A. Ta có $(d_2-d_1)/\lambda=6/2=3$, nên $M$ thuộc cực đại bậc 3.
}
\end{ex}

% ===== Câu 164 =====
% ID: L11C2B12-61-TN
% Mức: TH
\begin{ex}
% ID: L11C2B12-61-TN
% Mức: TH
Trong dạng “Xác định khoảng cách giữa các vân giao thoa”, Hai nguồn cùng pha, $\lambda=4$ cm. Điểm $N$ có $d_2-d_1=6$ cm là
\choice
{cực đại}
{\True cực tiểu}
{điểm đứng yên do không có sóng tới}
{không xác định}
\loigiai{
Đáp án B. $6/4=1,5$, là nửa nguyên nên $N$ thuộc cực tiểu.
}
\end{ex}

% ===== Câu 165 =====
% ID: L11C2B12-62-TN
% Mức: VD
\begin{ex}
% ID: L11C2B12-62-TN
% Mức: VD
Trong dạng “Xác định khoảng cách giữa các vân giao thoa”, Trên đoạn nối hai nguồn cùng pha, khoảng cách giữa hai cực đại liên tiếp bằng
\choice
{\True $\lambda/2$}
{$\lambda$}
{$2\lambda$}
{$\lambda/4$}
\loigiai{
Đáp án A. Trên đoạn nối hai nguồn, hiệu đường đi đổi hai lần theo li độ nên hai cực đại liên tiếp cách nhau $\lambda/2$.
}
\end{ex}

% ===== Câu 166 =====
% ID: L11C2B12-63-TN
% Mức: VD
\begin{ex}
% ID: L11C2B12-63-TN
% Mức: VD
Trong dạng “Xác định khoảng cách giữa các vân giao thoa”, Hiện tượng giao thoa là bằng chứng rõ ràng của
\choice
{\True tính chất sóng}
{tính chất hạt}
{sự bảo toàn khối lượng}
{chuyển động thẳng đều}
\loigiai{
Đáp án A. Giao thoa là hiện tượng đặc trưng của sóng.
}
\end{ex}

% ===== Câu 167 =====
% ID: L11C2B12-64-TN
% Mức: VD
\begin{ex}
% ID: L11C2B12-64-TN
% Mức: VD
Trong dạng “Xác định khoảng cách giữa các vân giao thoa”, Hai nguồn cùng pha cách nhau $10\,\text{cm}$, $\lambda=2$ cm. Số cực đại trên đoạn nối hai nguồn, kể cả hai nguồn nếu thỏa điều kiện, là
\choice
{9}
{10}
{\True 11}
{12}
\loigiai{
Đáp án C. Các bậc nguyên thỏa $|k|\le 5$, nên có 11 giá trị từ -5 đến 5.
}
\end{ex}

% ===== Câu 168 =====
% ID: L11C2B12-65-TN
% Mức: VDC
\begin{ex}
% ID: L11C2B12-65-TN
% Mức: VDC
Trong dạng “Xác định khoảng cách giữa các vân giao thoa”, Tại điểm có cực đại giao thoa, hai sóng thành phần đến đó
\choice
{\True cùng pha}
{ngược pha}
{vuông pha}
{khác tần số}
\loigiai{
Đáp án A. Cực đại xuất hiện khi hai dao động thành phần tăng cường nhau, tức cùng pha.
}
\end{ex}

% ===== Câu 169 =====
% ID: L11C2B12-66-TN
% Mức: VDC
\begin{ex}
% ID: L11C2B12-66-TN
% Mức: VDC
Trong dạng “Xác định khoảng cách giữa các vân giao thoa”, Nếu tần số nguồn tăng gấp đôi còn tốc độ truyền sóng không đổi thì bước sóng
\choice
{\True giảm một nửa}
{tăng gấp đôi}
{không đổi}
{tăng bốn lần}
\loigiai{
Đáp án A. Từ $\lambda=v/f$, khi $f$ tăng gấp đôi thì $\lambda$ giảm một nửa.
}
\end{ex}
